\documentclass[
ngerman,
paper=a4,
asymmetric,
fontsize=8pt,
% headings=small,
draft=false,
openany,
% parskip=half,
numbers=noenddot,
% listof=totoc,         % verschiebt Minitoc/Hyperref
% bibliography=totoc,   % verschiebt Minitoc/Hyperref
bibtotocnumbered,
% index=totoc,
% noonelinecaption,
footnotes=multiple,
toc=flat,
]{scrbook}
% \setlength{\parindent}{1.5em}
% \usepackage{scrhack}

\usepackage{luatextra}
\usepackage{fontspec}

% Optionen für Bilder:
% [draft] .. Es werden nur Rahmen angezeigt
% [final] .. Es werden Bilder normal angezeigt
% [demo} .. keine Bilder
\usepackage{graphicx} % Bilder-Support


%%%%%%%%%%%%%%%%%%%%%%%%%%%%%%%%%%%%%%%%%%%%%%%%%%%%%%%%%%%%%%%%%%%%%%%%%%%%%%%
%%%%%%%%% OPTIONEN
%%%%%%%%%%%%%%%%%%%%%%%%%%%%%%%%%%%%%%%%%%%%%%%%%%%%%%%%%%%%%%%%%%%%%%%%%%%%%%%
% optional --- Optionaler Text. 
\usepackage
[
%%% Status
% STATUS-DRAFT,
% STATUS-FINAL,
STATUS-EXPERIMENTAL,
% STATUS-DEV,
% STATUS-REVIEW,
% STATUS-FINAL,
%%% Git
%SHOW-GIT-COMMIT-ID,
%%%
% SHOW-TODO,
%%% Organisation
ORGASBFLO,
ORG-ASBFLO,
% ORG-ASBOE,
%%% Varianten
VAR-STANDARD,
VAR-RETTUNGSDIENST,
VAR-VAR0, % Stammvariante
% VAR-VAR1, % AASS, ASB Floridsdorf-Donaustadt
% VAR-VAR2, % ASBOE
%%% Region
REGAUT,
REG-AUT,
% REGGER;
REGVIE,
% REGNOE,   
]{optional}		% optional --- Optionaler Text. 
\input{./AASS-KOMP-2014.preamble.tex}

\usepackage[backend=bibtex8,
babel=other,
defernumbers=true,
natbib,
sorting=none,
sortlocale=de,
]{biblatex}

% 	\bibliography{AASS.bib}
    \addbibresource{BiblioDB-AASS.bib}
    \addbibresource{BiblioDB-ASBOe.bib}
	\addbibresource{BiblioDB-Med-Books.bib}
	\addbibresource{BiblioDB-Med-Books-Harrisons18De.bib}
	\addbibresource{BiblioDB-Med-Journals-Rettungsdienst.bib}
	\addbibresource{BiblioDB-Med-Links.bib}
	\addbibresource{BiblioDB-Med-Misc.bib}
	\renewcommand{\labelnamepunct}{\addcolon\space}
	\DeclareFieldFormat{title}{\mkbibbold{\mkbibemph{#1}}}
	\DeclareFieldFormat[article]{title}{\mkbibbold{\mkbibemph{#1}}}
	\DeclareFieldFormat[inbook]{title}{\mkbibbold{\mkbibemph{#1}}}
	% \DeclareFieldFormat{citetitle}{\mkbibbold{#1}}
	% 
	% \DeclareFieldFormat{citetitle}[article]{\mkbibbold{#1}}
	\DeclareFieldFormat{journaltitle}{\textsc{#1}}
	\DeclareFieldFormat[article]{journaltitle}{\textsc{#1}}
	\renewcommand{\bibfont}{\normalfont\tiny}
	\DeclareBibliographyCategory{PAM}
	\DeclareBibliographyCategory{BestOf}
	\defbibnote{bestof}{Dieser Abschnitt beinhaltet besonders
	lesenswerte Literaturreferenzen. Eine vollständgie
	Auflistung der Quellen findet sich im Folgeabschnitt.
	Sofern dieser nicht Teil Ihrer Distribution ist, finden
	Sie sie auch online unter
	\url{http://www.aass.at/download}.}
% \usepackage{asboe}

% Trennungen
\input{./AASS.hyphenation.tex}

% Dokumenteinfos (Status, Version, ...)
% \RequirePackage[iso]{isodateo}
\RequirePackage{datetime}
\renewcommand{\dateseparator}{-}
\newdateformat{iso}{\THEYEAR-\twodigit{\THEMONTH}-\twodigit{\THEDAY}}
\newtimeformat{nummeric}{\twodigit{\THEHOUR}\twodigit{\THEMINUTE}}
\settimeformat{nummeric}


\newcommand{\VarAassVersion}{2.0}
\newcommand{\DocVersion}{2014.0-dev-\pdfdate}
% \renewcommand{\DocVersion}{0.20}
\newcommand{\DocVersionTyp}{Entwicklungsversion}
% \renewcommand{\DocVersionTyp}{Vorabversion}
% \renewcommand{\DocVersionTyp}{Release}
\newcommand{\DocVersionInfo}{Entwicklungsversion -- Für den internen Gebrauch.}
% \renewcommand{\DocVersionInfo}{Vorabversion -- F.\,d. int. Gebrauch.} 
\renewcommand{\DocVersionInfo}{Release} 
\hypersetup{
% draft,
unicode=true,
colorlinks=false,
linkcolor=darkBlue,
citecolor=darkBlue,
filecolor=darkBlue,
pagecolor=darkBlue,
urlcolor=darkBlue,
% Rahmen um Links:
pdfborderstyle={/S/U/W 0},  % Unterstreichung
% pdfborderstyle={/S/U/W 0.5},  % Unterstreichung
pdfborder={0 0 0},          % Dicke
% linkbordercolor=DefaultB,     % Farbe
% urlbordercolor=DefaultB,     % Farbe
pdftitle={Arbeits- und Ausbildungsstandards für den Sanitätsdienst \DocVersion},
% Sets the document information Title field pdfauthor=ARGE AASS, % Sets the document information Author field
pdfsubject={AASS \DocVersion}, % Sets the document information Subject field
pdfcreator=ARGE AASS, % Sets the document information Creator field
pdfproducer=ARGE AASS, % Sets the document information Producer field
pdfkeywords=ARGE AASS, % Sets the document information Keywords field
% pdfpagelayout=TwoPageRight,
}
% By default, URLs are printed using mono-spaced fonts. If you
% don't like it and you want them to be printed with the same style 
% of the rest of the text, you can use this:
\urlstyle{same}
% Staus Möglichkeiten:
% Final 				Endversion
% DraftLight			+ Bilder, - Kommentare
% DraftMedium		- Bilder, - Kommentare
% DraftFull			- Bilder, + Kommenatre
% DraftSuper			- Bilder, + Kommenatare, + draft=yes
\newcommand{\DocStatus}{DraftFull}
\newcommand{\DocGitCommit}{Unofficial build, no SHA1 included!}
% \renewcommand{\DocGitCommit}{\input{../tmp/git-lastcommitsha1.tmp}}
\newcommand{\DocGitCommitFormatted}{\ttfamily\DocGitCommit}
\newcommand{\DocGitCommitTiny}{\tiny\DocGitCommitFormatted}
\newcommand{\DocGitCommitScriptsize}{\scriptsize\DocGitCommitFormatted}
\newcommand{\DocGitCommitFootnotesize}{\footnotesize\DocGitCommitFormatted}

% Kopf- und Fusszeilen
%%%%%%%%%%%%%%%%%%%%%%%%%%%%%%%%%%%%%%%%%%%%%%%%%%%%%%%%%%%
%%% Kopf- und Fußzeilen
\usepackage{fancyhdr}
\pagestyle{fancy}

\newcommand{\Fancyheadreset}{
  \fancyheadoffset[RE,RO]{\marginparsep+\marginparwidth}
  \fancyfootoffset[RE,RO]{\marginparsep+\marginparwidth}
  }


\pagestyle{fancy}
\fancyhf{} % clear all header and footer fields
\fancyheadoffset[RE,RO]{\marginparsep+\marginparwidth}
\fancyfootoffset[RE,RO]{\marginparsep+\marginparwidth}

\renewcommand{\chaptermark}[1]{\markboth{#1}{}}
\renewcommand{\sectionmark}[1]{\markright{\thesection.\ #1}}
\renewcommand{\headrule}{{\color{Grey}%
\hrule width\headwidth height\headrulewidth \vskip-\headrulewidth}}












% \usepackage[scaled]{helvet}

% % TikZ-Pakete
% \usepackage{tikz}
% \usetikzlibrary{positioning,backgrounds,shapes}
% 
% % Kopf- und Fußzeile
% \usepackage[automark]{scrpage2}
% \clearscrheadings
% \clearscrplain
% 
% \tikzset{pagenumber/.style={rectangle,rounded corners=5pt,inner sep=6pt, fill=gray!20,draw=gray!30,line width=0.1em}}
% 
% \lehead{\begin{tikzpicture}[remember picture,overlay]
% \begin{pgfonlayer}{background}
% \draw[fill=gray!10, fill opacity=1.0, draw=none] (-.65\paperwidth,0cm) rectangle (\paperwidth+2cm,2cm);
% \end{pgfonlayer}
% \path (current page.north west) node[anchor=west,xshift=1.58cm, yshift=-0.9cm,pagenumber] {\parbox{\widthof{\Large 11111}}{\centering \color{black} \bfseries\Large \thepage}};
% \end{tikzpicture}}
% 
% \rehead{\begin{tikzpicture}[remember picture,overlay]
% \path (current page.north east) node[anchor=east,xshift=-3.4cm, yshift=-0.9cm] {\color{black}\upshape \leftmark};
% \end{tikzpicture}}
% 
% \lohead{\begin{tikzpicture}[remember picture,overlay]
% \begin{pgfonlayer}{background}
% \draw[fill=gray!10, fill opacity=1.0, draw=none] (-.65\paperwidth,0cm) rectangle (\paperwidth+2cm,2cm);
% \end{pgfonlayer}
% \path (current page.north west) node[anchor=west,xshift=3.4cm, yshift=-1.0cm] {\color{black}\upshape \headmark};
% \end{tikzpicture}}
% 
% \rohead{\begin{tikzpicture}[remember picture,overlay]
% \path (current page.north east) node[anchor=east,xshift=-1.58cm, yshift=-0.9cm,pagenumber] {\parbox{\widthof{\Large 11111}}{\centering \color{black}\bfseries\Large \thepage}};
% \end{tikzpicture}}
% \setlength{\headsep}{1.1cm}
% 
% \lefoot{{\noindent\AASS{}~ \sffamily\tiny v.\,\DocVersion~~|~~\url{http://www.aass.at}}}
% \refoot{{\noindent\DocGitCommitTiny}}
% 
% \lofoot{{\noindent\DocGitCommitTiny}}
% \rofoot{{\noindent{\sffamily\tiny\url{http://www.aass.at}~~|~~v.\,\DocVersion}~ \AASS{}}}
















%%% HEADER %%%
% right
\fancyhead[RO]{{\sffamily{\color{BaseLightVery}\CC{[\ChapterCode]}~\rightmark~|~}{\color{Base}\thepage}}}
\fancyhead[RO]{\begin{tikzpicture}[overlay,remember picture]
      \fill [color=AltBlueLightSuper]
        ($(current page.north east) - (0pt,0pt)$)
        rectangle
        ($ (current page.south east) - (138.5pt,0cm) $);
  \end{tikzpicture}{\sffamily{\color{BaseLightVery}\CC{[\ChapterCode]}~\rightmark~|~}{\color{Base}\thepage}}
}
% \fancyhead[RE]{\EE{\protect\AuthNotes{ | {\protect\today} | ENTWURF |}}}
% left
% \fancyhead[LO]{\EE{\protect\AuthNotes{ | ENTWURF | {\protect\today} |}}}
\fancyhead[LE]{\begin{tikzpicture}[overlay,remember picture]
      \fill [color=AltBlueLightSuper]
        ($(current page.north east) - (0pt,0pt)$)
        rectangle
        ($ (current page.south east) - (161.60144pt,0cm) $);
  \end{tikzpicture}{\sffamily{\color{Base}\thepage}{\color{BaseLightVery}~|~\leftmark}}}
% %%% DRAFT %%%
% \fancyhead[RE,LO]{{\noindent\color{red}\ttfamily\footnotesize\upshape{\textbf{\DocVersionTyp}} | \textbf{Nur für den internen Gebrauch!}}}

%%% FOOTER %%%
% center
\fancyfoot[CE,CO]{{{\sffamily\tiny  \DocVersionInfo}}}
% odd
\fancyfoot[LO,RE]{{\noindent\DocGitCommitTiny}}
\fancyfoot[LO,RE]{{\noindent\tiny\FontSchemeCode}}
\fancyfoot[RO]{{\noindent{\sffamily\tiny\url{http://www.aass.at}~~|~~\DocVersion}~ \AASS{}}}
% even
\fancyfoot[LE]{{\noindent\AASS{}~ \sffamily\tiny \DocVersion~~|~~\url{http://www.aass.at}}}

% %%% OVERWRITE SETTINGS FOR DRAFT %%%
% \fancyfoot[CE,CO]{{\noindent\color{red}\ttfamily\footnotesize\upshape{\textbf{\DocVersionTyp}} | \textbf{Nur für den internen Gebrauch!}}}

%

\opt{VAR-VAR2}{
% \fancyhf{} % clear all header and footer fields
% center
\fancyfoot[CE,CO]{{{\noindent\sffamily\tiny\color{AsboeGreyLight}\DocVersionInfo~~|~~{\noindent\DocGitCommitTiny}}}}
% odd
\fancyfoot[LO]{{\noindent\sffamily\tiny\color{AsboeGreyLight}\url{http://www.aass.at}~~|~~\url{http://www.samariterbund.net}}}
\fancyfoot[RO]{{\noindent{\sffamily\tiny\color{AsboeGreyLight}\DocVersion}~~\AASS{}\,{\color{AsboeGreyLight}|}\,{\color{AsboeRed}ASBÖ}~ }}
% even
\fancyfoot[LE]{{\noindent\AASS{}\,{\color{AsboeGreyLight}|}\,{\color{AsboeRed}ASBÖ}~~\sffamily\tiny\color{AsboeGreyLight}\DocVersion}}
\fancyfoot[RE]{{\noindent\sffamily\tiny\color{AsboeGreyLight}\url{http://www.aass.at}~~|~~\url{http://www.samariterbund.net}}}
}


% % %



% Redefining pagestyle plain
\fancypagestyle{plain}{%
\fancyhf{} % clear all header and footer fields
%%% IMMER %%%
% center
\fancyfoot[CE,CO]{{{\sffamily\tiny  \DocVersionInfo}}}
% odd
\fancyfoot[LO]{{\noindent\DocGitCommitTiny}}
\fancyfoot[RO]{{\noindent{\sffamily\tiny\url{http://www.aass.at}~~|~~\DocVersion}~ \AASS{}}}
% even
\fancyfoot[LE]{{\noindent\AASS{}~ \sffamily\tiny \DocVersion~~|~~\url{http://www.aass.at}}}
\fancyfoot[RE]{{\noindent\DocGitCommitTiny}}

\opt{VAR-VAR2}{
\fancyhf{} % clear all header and footer fields
% center
\fancyfoot[CE,CO]{{{\sffamily\tiny\color{AsboeGreyLight}\DocVersionInfo~|~{\noindent\DocGitCommitTiny}}}}
% odd
\fancyfoot[LO]{{\noindent\DocGitCommitTiny}}
\fancyfoot[RO]{{\noindent{\sffamily\tiny\url{http://www.aass.at}~~|~~\DocVersion}~ \AASS{}}}
% even
\fancyfoot[LE]{{\noindent\AASS{}\,{\color{AsboeGreyLight}|}\,{\color{AsboeRed}ASBÖ}~ \sffamily\tiny \DocVersion}}
\fancyfoot[RE]{{\color{AsboeGreyLight}\url{http://www.aass.at}~~|~~\url{http://www.samariterbund.net}}}
}

% \fancyfoot[CE,CO]{{{\sffamily\tiny {\DocVersionTyp} | Nur für den internen Gebrauch!}}}
%%% DRAFT %%%
% \fancyfoot[CE,CO]{{\noindent\color{red}\ttfamily\footnotesize\upshape{\textbf{\DocVersionTyp}} | \protect\today{} | \textbf{Nur für den internen Gebrauch!}}}
%
% % {\tiny
% % Zweig: }\input{./tmp/git-current-branch.tmp} -- {\tiny
% % Letztes Commit: }\input{./tmp/git-lastcommit.tmp}
% }}
%
% %%% FINAL %%%

\renewcommand{\headrulewidth}{0pt}
\renewcommand{\footrulewidth}{0pt}
}


% Ausgeschlossene Kapitel
% \usepackage{excludeonly}
% %\excludeonly{include-betriebsintern-chapter,include-regional-chapter,include-taetigkeitsbild-chapter,include-erstehilfe-chapter,include-recht-chapter,include-reserve-ende}
\excludeonly{
./chapters/chapter-NEW,
./chapters/chapter-ARZ,
./chapters/chapter-WIS,
% ./chapters/chapter-INT,
}

% Schliesse nur bestimmte Kapitel ein ...
% % \includeonly{%
% ./chapters/chapter-ABD,
% ./chapters/chapter-AED,
% ./chapters/chapter-ALS,
% ./chapters/chapter-ANA,
% ./chapters/chapter-ARZ,
% ./chapters/chapter-ASB,
% ./chapters/chapter-ASS,
% ./chapters/chapter-AUP,
% ./chapters/chapter-BAU,
% ./chapters/chapter-BET,
% ./chapters/chapter-BLS,
% ./chapters/chapter-DOK,
% ./chapters/chapter-EHI,
% ./chapters/chapter-EKG,
% ./chapters/chapter-EIN,
% ./chapters/chapter-GEF,
% ./chapters/chapter-GLO,
% ./chapters/chapter-GYN,
% ./chapters/chapter-HKL,
% ./chapters/chapter-HYG,
% ./chapters/chapter-INK,
% ./chapters/chapter-INT,
% ./chapters/chapter-JUF,
% ./chapters/chapter-JUS,
% ./chapters/chapter-KHD,
% ./chapters/chapter-KIN,
% ./chapters/chapter-KRI,
% ./chapters/chapter-LER,
% ./chapters/chapter-MP1,
% ./chapters/chapter-NEU,
% ./chapters/chapter-NEW,
% ./chapters/chapter-NKA,
% ./chapters/chapter-NWG,
% ./chapters/chapter-ADVANCED-Atemweg,
% ./chapters/chapter-PAM,
% ./chapters/chapter-PKU,
% ./chapters/chapter-PSI,
% ./chapters/chapter-PSY,
% ./chapters/chapter-REG,
% ./chapters/chapter-RSP,
% ./chapters/chapter-RWT,
% ./chapters/chapter-SAN,
% ./chapters/chapter-SHF,
% ./chapters/chapter-SPE,
% ./chapters/chapter-STW,
% ./chapters/chapter-SYM,
% ./chapters/chapter-TAE,
% ./chapters/chapter-THE,
% ./chapters/chapter-TOD,
% ./chapters/chapter-TOX,
% ./chapters/chapter-TRA,
% ./chapters/chapter-URO,
% ./chapters/chapter-VIT,
% ./chapters/chapter-WEI,
% ./chapters/chapter-WIS,
% ./chapters/chapter-WUN,
% inc-APP,
% inc-RES-ende,
% ./chapters/part-Lehrmaterialien,
% }



\setcounter{errorcontextlines}{10}



\renewcommand{\Abk}[1]{#1}



% \renewcommand{\MassnahmenNG}[4]
% 	{%
% % 	#3
% % 	\par
% 	\begin{multicols}{2}[\protect\section{{\color{Indigo}#1:} \E{#2}}\index{#2}#3]
% \CorrectionListsBox#4
% 	\end{multicols}
% 	}

\renewcommand{\MassnahmenNG}[4]
	{%
% 	#3
% 	\par
	\section{{\color{Indigo}#1:} \E{#2}}\index{#2}#3
{\CorrectionListsBox#4}
	}

\renewcommand{\MassnahmenWide}[4]
	{%
% 	#3
% 	\par
	\section{{\color{Indigo}#1:} \E{#2}}\index{#2}#3
{\CorrectionListsBox#4}
	}
	
\usepackage{AassStyleCollectionProcedures}
\usepackage{AassStyleCollectionCoverFront}

% \usepackage{xr}
% 	\opt{VAR-VAR1}{\externaldocument[AASS-]{AASS-0.26-var1-Archiv}}
% 	\opt{VAR-VAR2}{\externaldocument[AASS-]{AASS-0.26-var2-Archiv}}
% \externaldocument[AASS-]{AASS}

% \renewcommand{\MassnahmenNG}[4]{#1
% 
% #2
% 
% #3
% 
% #4}

\def\ChapterCode{---}


% #######################################################################################################
% #######################################################################################################
% #######################################################################################################
% #######################################################################################################
% #######################################################################################################
% #######################################################################################################
% #######################################################################################################
% BEGIN DOCUMENT ########################################################################################
% #######################################################################################################
\begin{document}
\MultiColListStart{toc}{2}
\MultiColListStart{lofazit}{2}
\MultiColListStart{lomassnahmen}{2}
\DTLloaddb[headers={CODE,VERSION,UPDATE,TYPE,TITLE,REMARKS,TARGET,SUMMARY,SPECIALTY,ICD10-P,ICD10-1, ICD10-2,ICD10-3,ICD10-4,MODIFIER,LEVEL,ICD10,ICD10-TEXT}]{procedures}{procedures-export.csv}



\newgeometry{verbose,
a5paper,
asymmetric,
reversemarginpar,
rmargin=13mm,
lmargin=13mm,
top    = 10mm,
bottom = 10mm,
marginparwidth = 0mm,
footskip = 5mm,
bindingoffset=1cm}

\savegeometry{FullPage}

\newgeometry{verbose,
a5paper,
asymmetric,
reversemarginpar,
rmargin=7mm,
lmargin=7mm,
top    =7mm,
bottom =7mm,
marginparwidth = 0mm,
footskip = 5mm,
bindingoffset=1.0cm}
\savegeometry{WholePage}
\restoregeometry



\opt{STATUS-DRAFT-SHOWKEYS}{\renewcommand{\label}[1]{{\color{black}\fbox{{\ttfamily #1 }}}}}


\loadgeometry{FullPage}

% 1: Title
% 2: Subtitle
% 3: Head
% 4: Picture
% 5: Editors
% 6: Bottom
% 7: Publisher
\CoverFrontMake
{Arbeits- und Ausbildungsstandards\\für den Sanitätsdienst}
{Kompendium}
{Head}
{Picture}
{Editors}
{Bottom}
{\ARGE}

\clearpage
\thispagestyle{empty}


%%% uppertitleback

	Covergestaltung: \emph{Sebastian Gabriel.}

    {Titelbild: \emph{Christoph Pallinger} }
    
\vfill
    
\begin{WidePar}\centering\begin{minipage}{8cm}
    \begin{center} 
        Dieses Projekt entstand in Zusammenarbeit mit dem
        
        \emph{\bfseries Ausbildungszentrum des ASB Floridsdorf-Donaustadt} 
        
        und wurde maßgeblich von der
        
        \emph{\bfseries Cyberservice Internetdienstleistungsgesellschaft~m.b.H.}
        
        und
        
        \emph{\bfseries topfen.net} 
        
        unterstützt.\vspace{2em}
    
        \includegraphics[width=3cm]{./icons/asb.png}
        
        {\sffamily\fontfamily{phv}Gruppe Floridsdorf-Donaustadt}\par\vspace{2em}
        
        \includegraphics[width=5cm]{./images/logos/cyberservice.png}
    
        {\sffamily\fontfamily{phv}www.cyberservice.net}\par\vspace{2em}
    
        {\Huge topfen.net}\par\vspace{1em}
        \textit{serving the community} \par \textit{since 2001}
    \end{center}
    \end{minipage}
    \vspace{2em}
    \begin{center}
%         {\bfseries\itshape Nur für den internen Gebrauch gem §\,42 Abs.\,6
%         UrheberG bestimmt. Weitergabe nicht gestattet.}
    \end{center}
    
    \end{WidePar}

\clearpage
\thispagestyle{empty}

~
\vspace{3cm}

\begin{center}
saluti et solatio aegrorum
\end{center}



\bigskip

\bigskip

\bigskip

\bigskip

\bigskip

\bigskip

\bigskip

\bigskip

\bigskip

\bigskip

\dictum[Andrew S. Tanenbaum (übersetzt)]{Das Schöne an Standards
ist, es gibt so viele, aus denen man auswählen kann.}

\vfill

\begin{center}
\Speech{"`Früher war mehr Lametta."'}

\bigskip

\bigskip

in memoriam

Vicco von Bülow (\VonBis{1923}{2011})


\vspace{3cm}

\end{center}





































\opt{STATUS-DRAFT}{{\scriptsize\listof{warning}{Verzeichnis der Warnungen}}}



\restoregeometry

\addsec{Benutzungshinweise}

Der vorliegende Maßnahmenkatalog ist ein Begleitwerk zu den
Arbeits- und Ausbildungsstandards für den Sanitätsdienst
(\AASS). Die zum Zeitpunkt des Redaktionsschlusses
korrespondierende Ausgabe ist die Version \VarAassVersion.
Alle Verweise auf die {\AASS} 
beziehen sich auf diese Version, und können von einer
späteren Ausgabe abweichen.

\addsec{Lizenz der {\AASS}}

Die Inhalte der {\AASS} als Gesamt- und Sammelwerk unterliegen 
den Nutzungsbedingungen der Creative-Commons-Lizenz \Speech{\enquote*{Namensnennung --
Keine Bearbeitungen 3.0 Österreich}} (\I{CC-BY-ND/3.0/at}, 
abrufbar
unter
\url{http://creativecommons.org/licenses/by-nd/3.0/at/}).

Davon nicht berührt sind die Lizenzbedingungen von eingearbeiteten Werken
(z.\,B. Bilder und Grafiken), die davon abweichende
Lizenzbestimmungen aufweisen können.

{\minisec{Kurzfassung der Nutzungsbestimmungen (unverbindlich):}

\begin{itemize}
  \item Sie dürfen dieses Werk
  \begin{itemize}
    \item das Werk bzw. den Inhalt vervielfältigen, verbreiten und öffentlich zugänglich machen
    \item das Werk kommerziell nutzen
  \end{itemize}
  \item Zu den folgenden Bedingungen: 
\begin{itemize}
  \item \EE{Namensnennung} -- Sie müssen den Namen des
  Autors/Rechteinhabers in der von ihm festgelegten Weise
  nennen.
    \item \EE{Keine Bearbeitung} -- Dieses Werk bzw. dieser
    Inhalt darf nicht bearbeitet, abgewandelt oder in
    anderer Weise verändert werden.
  \end{itemize}
\end{itemize}
\EE{Hinweis} -- Im Falle einer Verbreitung müssen Sie
anderen alle Lizenzbedingungen mitteilen, die für dieses
Werk gelten. Am einfachsten ist es, an entsprechender Stelle
einen Link auf \url{http://creativecommons.org/licenses/by-nd/3.0/at/} einzubinden.
}


%\renewcommand\contentsname{Inhalt - Schnell\"ubersicht}
\dominitoc[n]

\begin{WidePar}
%   {\setcounter{tocdepth}{3}\scriptsize\tableofcontents}
% \tableofcontents
\end{WidePar}
%%%%%%%%%%%%%%%%%%%%%%%%%%%%%%%%%%%%%%%%%%%%%%%%%%%%%%%%%%%
%%% Index
% \clearpage
{\small\printindex}


\chapter{Erste-Hilfe}




% % M %%%%%%%%%%%%%%%%%%%%%%%%%%%%%%%%%%%%%%%%%%%%%%%%%%% M %
% % 
% \ProcedurePrintDefault{}
\MassnahmenNG{Erste-Hilfe-Maßnahmen}{Bewusstlosigkeit}{\label{m:ehi-bewusstlosigkeit}}{
\begin{itemize}
    \item Stabile Seitenlage
\end{itemize}
}




% % M %%%%%%%%%%%%%%%%%%%%%%%%%%%%%%%%%%%%%%%%%%%%%%%%%%% M %
% % 
% \ProcedurePrintDefault{}
\MassnahmenNG{Erste-Hilfe-Maßnahmen}{Atem- und
Kreislaufstillstand}{\label{m:ehi-atem-kreislaufstillstand}}{
\begin{itemize}
  \item Herz-Lungen-Wiederbelebung (Reanimation)
  \begin{itemize}
    \item Herzdruckmassage
    \item Beatmung
  \end{itemize}
\end{itemize}
}




% % M %%%%%%%%%%%%%%%%%%%%%%%%%%%%%%%%%%%%%%%%%%%%%%%%%%% M %
% % 
% \ProcedurePrintDefault{}
\MassnahmenNG{Erste-Hilfe-Maßnahmen}{Starke
Blutung}{\label{m:ehi-starke-blutung}}{
\begin{itemize}
  \item Blutstillung
  \begin{itemize}
    \item Hochhalten der betroffenen Gliedmaße
    \item Zudrücken,
    \item Abdrücken einer blutzuführenden Arterie,
    \item Druckverband,
    \item Abbindung.
  \end{itemize}
\end{itemize}

}




% % M %%%%%%%%%%%%%%%%%%%%%%%%%%%%%%%%%%%%%%%%%%%%%%%%%%% M %
% % 
% \ProcedurePrintDefault{}
\MassnahmenNG{Erste-Hilfe-Maßnahmen}{Schockbekämpfung}{\label{m:ehi-schockbekaempfung}}{

\begin{itemize}
  \item Ursache beseitigen (wenn möglich), ggfs. Blutstillung
  \item Beengende Kleidungsstücke öffnen
  \item Situationsgerechte Lagerung
  \item Wärmeerhalt
  \item Beruhigender Zuspruch
  \item Beobachtung
\end{itemize}
}




% % M %%%%%%%%%%%%%%%%%%%%%%%%%%%%%%%%%%%%%%%%%%%%%%%%%%% M %
% % 
% \ProcedurePrintDefault{}
\MassnahmenNG{Erste-Hilfe-Maßnahmen}{Bagatellverletzungen}{\label{m:ehi-bagatellverletzungen}}{

\begin{itemize}
  \item Wunde reinigen
  \item Evtl. steriler Wundverband
  \item Auf den Tetanusschutz achten
\end{itemize}
}




% % M %%%%%%%%%%%%%%%%%%%%%%%%%%%%%%%%%%%%%%%%%%%%%%%%%%% M %
% % 
% \ProcedurePrintDefault{}
\MassnahmenNG{Erste-Hilfe-Maßnahmen}{Verbrennungen}{\label{m:verbrennung-eh}}{
\begin{itemize}
  \item Hitzeeinwirkung sofort stoppen.
  \item Wenn 
  \begin{itemize}
    \item innerhalb von 2 Minuten
  möglich,
    \item Ausdehnung <\,20\PC* (Kinder <\,10\PC, Säuglinge
    <\,10\PC*) 
  \end{itemize}  für \E{maximal 10\Min* mit} mit
    handwarmen Wasser kühlen (Richtwert 20\,\DegC); 
    
    andernfalls ist keine Kühlung durchzuführen
  \cite{Asanger2010,ASBOEAkademieUpdate2010Verbrennung}.
  
  \Ca{Achtung: Gefahr der Unterkühlung!}
  \item Wärmeerhalt ist besonders wichtig!
  \item Keimfreier Wundverband
  \item Schockbekämpfung
  \item \Ca{Keine Salben oder Puder etc. verwenden}
\end{itemize}
}




% % M %%%%%%%%%%%%%%%%%%%%%%%%%%%%%%%%%%%%%%%%%%%%%%%%%%% M %
% % 
% \ProcedurePrintDefault{}
\MassnahmenNG{Erste-Hilfe-Maßnahmen}{Erfrierungen}{\label{m:ehi-erfrierungen}}{

\begin{itemize}
  \item Keimfreier Wundverband
  \item Restlichen Körper aufwärmen
  \item Warme, gezuckerte Getränke verabreichen
  \item Keinen Alkohol zu trinken geben!
  \item Nicht frottieren! 
%   % Erfrierungen sind lokal, daher hier kein Hinweis auf
%   % Bergungstod
%   sog.
%   \Ca{Bergungstod!}\footnote{Unter dem Begriff \emph{Bergungstod} versteht
%   man das Versterben des Patienten während der Rettung in Folge eines Blutrückflusses
%   aus der kalten Körperschale in den wärmeren Körperkern, welcher durch die
%   Bewegung des Patienten während der Rettung ausgelöst wird.}
\end{itemize}
}




% % M %%%%%%%%%%%%%%%%%%%%%%%%%%%%%%%%%%%%%%%%%%%%%%%%%%% M %
% % 
% \ProcedurePrintDefault{}
\MassnahmenNG{Erste-Hilfe-Maßnahmen}{Unterkühlung}{\label{m:ehi-unterkuehlung}}{

\begin{itemize}
  \item Aufwecken und wach halten
  \item Vorsichtig in Decken oder Kleidungsstücke einwickeln
  \item Rettungsdienst verständigen
  \item Warme Getränke verabreichen, wenn Patient bewusstseinsklar ist
  \item Keinen Alkohol zu trinken geben!
  \item Keinesfalls bewegen oder frottieren: sog.
  \Ca{Bergungstod!}\footnote{Unter dem Begriff \emph{Bergungstod} versteht
  man das Versterben des Patienten während der Rettung in Folge eines Blutrückflusses
  aus der kalten Körperschale in den wärmeren Körperkern, welcher durch die
  Bewegung des Patienten während der Rettung ausgelöst wird.}
\end{itemize}
}




% % M %%%%%%%%%%%%%%%%%%%%%%%%%%%%%%%%%%%%%%%%%%%%%%%%%%% M %
% % 
% \ProcedurePrintDefault{}
\MassnahmenNG{Erste-Hilfe-Maßnahmen}{Verstauchung}{\label{m:ehi-verstauchung}}{

\begin{itemize}
  \item Ruhigstellung
  \item Kalte Umschläge
  \item Unfallabteilung aufsuchen
  \item Hochlagern % ZWING 2010-03-xx
\end{itemize}
}




% % M %%%%%%%%%%%%%%%%%%%%%%%%%%%%%%%%%%%%%%%%%%%%%%%%%%% M %
% % 
% \ProcedurePrintDefault{}
\MassnahmenNG{Erste-Hilfe-Maßnahmen}{Verrenkung}{\label{m:ehi-verrenkung}}{

\begin{itemize}
    \item Vorgefundene Stellung beibehalten
    \item Ruhigstellung
    \item Rettungsdienst verständigen
    \item Keine Einrenkversuche!
\end{itemize}
}




% % M %%%%%%%%%%%%%%%%%%%%%%%%%%%%%%%%%%%%%%%%%%%%%%%%%%% M %
% % 
% \ProcedurePrintDefault{}
\MassnahmenNG{Erste-Hilfe-Maßnahmen}{Knochenbrüche}{\label{m:ehi-knochenbrueche}}{

\begin{itemize}
    \item Ruhig liegen lassen
    \item Wundversorgung bei offenen Knochenbrüchen
    \item Beengende Kleidungsstücke lockern
    \item Schuhe öffnen, nicht ausziehen
    \item Ruhigstellung durch unterstützte Lagerung
    \item Schockbekämpfung
\end{itemize}
}





% % M %%%%%%%%%%%%%%%%%%%%%%%%%%%%%%%%%%%%%%%%%%%%%%%%%%% M %
% % 
% \ProcedurePrintDefault{}
\MassnahmenNG{Erste-Hilfe-Maßnahmen}{Vergiftungen}{\label{m:ehi-vergiftungen}}{
\begin{itemize}
    \item Patient nicht ansprechbar:
    \begin{itemize}
        \item Retten aus der Gefahr (Selbstschutz!)
        \item Lebensrettende Sofortmaßnahmen
    \end{itemize}
    \item Patient ansprechbar:
    \begin{itemize}
        \item Wenn das Gift bekannt ist, Vergiftungsinformationszentrale
        (\Code{01\,/\,406\,43\,43}) anrufen und \E{Expertenrat} einholen.
        \item Wenn Gift unbekannt: Betreuen und → Notruf/Krankenhaus.
        \item Selbstschutz!
        \item Giftstoffe sichern und mitnehmen.
    \end{itemize}
\end{itemize}
}




% % M %%%%%%%%%%%%%%%%%%%%%%%%%%%%%%%%%%%%%%%%%%%%%%%%%%% M %
% % 
% \ProcedurePrintDefault{}
\MassnahmenNG{Erste-Hilfe-Maßnahmen}{Verätzungen der
Haut}{\label{m:ehi-veraetzungen-haut}}{

\begin{itemize}
   \item \Ca{Selbstschutz} beachten!
    \item Mit ätzender Substanz getränkte Kleidung entfernen
    \item Haut mit reinem Wasser abspülen
    \item Keimfreier Wundverband
    \item Evtl. Schockbekämpfung
\end{itemize}
}




% % M %%%%%%%%%%%%%%%%%%%%%%%%%%%%%%%%%%%%%%%%%%%%%%%%%%% M %
% % 
% \ProcedurePrintDefault{}
\MassnahmenNG{Erste-Hilfe-Maßnahmen}{Verätzungen der
Augen}{\label{m:ehi-veraetzungen-augen}}{

\begin{itemize}
    \item Auge mit reinem Wasser von innen nach außen ausspülen. Abwasser
    \E{nicht über das gesunde Auge} abrinnen lassen
    \item Auge mit keimfreiem Wundverband bedecken; zum Ruhigstellen \E{beide}
    Augen bedecken
    \item Selbstschutz beachten!
\end{itemize}
}




% % M %%%%%%%%%%%%%%%%%%%%%%%%%%%%%%%%%%%%%%%%%%%%%%%%%%% M %
% % 
% \ProcedurePrintDefault{}
\MassnahmenNG{Erste-Hilfe-Maßnahmen}{Verätzungen des
Verdauungstraktes}{\label{m:ehi-veraetzungen-verdauungstrakt}}
{ 
\begin{itemize}
    \item Mund mit Wasser ausspülen
    \item Nicht erbrechen lassen!
    \item Substanz bekannt → Vergiftungsinformationszentrale anrufen
    \item evtl. Schockbekämpfung
    \item Selbstschutz beachten!
\end{itemize}
}
























\chapter{Allgemeine Sanitätshilfemaßnhamen}




% % M %%%%%%%%%%%%%%%%%%%%%%%%%%%%%%%%%%%%%%%%%%%%%%%%%%% M %
% % 
% \ProcedurePrintDefault{}
\MassnahmenNG{Spezielle Maßnahmen}{Sauerstoffberieselung}
{\label{m:sauerstoffberieselung}}
{%
Bei jedem Patienten, bei dem eine
lebensbedrohliche Störung einer vitalen Funktion eingetreten
ist oder einzutreten droht (\enquote{Notfallpatient}), soll,
sofern keine Kontraindiktaionen vorliegen, soviel Sauerstoff
verabreicht werden, sodass die Sauerstoffsättigung
(Sp\ce{O2}) im Bereich von \EE{\VonBis{94}{98}\PC*} erreicht wird.

\begin{enumerate}
  \item Kontraindikationen und Gegenanzeigen prüfen:
  \begin{itemize}
    \item COPD (\FullPrettyRef{AASS-topic:copd})
    \item Hyperventilationssyndrom, Hyperventilationstetanie
    (\FullPrettyRef{AASS-topic:hyperventailationssyndrom})
  \end{itemize}
  \item Situationsgerechte Dosierung je nach
  zugrundeliegender Erkrankung. Grundsätzlich soll ein
  Sp\ce{O2} von
\EE{\VonBis{94}{98}\PC*} erreicht werden. Steht keine Pulsoxymetrie zur Verfügung, ist als
Richtwert von einer Dosis von \EE{8\LPerMin*} auszugehen,
welche dem klinischen Zustand des Patienten angepasst werden
muss.
  \item Auswahl des Hilfsmittels:
  \begin{itemize}
    \item \E{Sauerstoff\-brille}
    \ce{O2}-Flow bis~5\LPerMin*,
    \item \E{Sauerstoff\-maske}
    \ce{O2}-Flow ab 5\LPerMin*,
    \item \E{Sauerstoff\-maske mit
    Reservoir} \ce{O2}-Flow ab 5\LPerMin*,
  \end{itemize}
  \item Patientenaufklärung
  \item Sauerstoffgerät einschalten und Durchflußrate
  einstellen
  \item Bei Verwendung einer Sauerstoff\-maske mit
  Reservoir: Reservoir füllen
  \item Sauerstoffbrille/-maske positionieren 
\end{enumerate} 

\cite{AsboeLehrmeinungRd2011-09}
}


\cleardoublepage

\chapter{Allgemeines Patientenmanagement}

\label{apm}





% % M %%%%%%%%%%%%%%%%%%%%%%%%%%%%%%%%%%%%%%%%%%%%%%%%%%% M %
% % 
% \ProcedurePrintDefault{}
\MassnahmenNG{Standardmaßnahmen}{Immer durchzuführende
Standardmaßnahmen}
{\label{m:standardmassnahmen-immer}}{%
Die folgenden Maßnahmen sind grundsätzlich immer in einer \E{der Situation
angemessenen} Art und Weise durchzuführen. Im begründeten Ausnahmefall kann es
allerdings notwendig oder sinnvoll sein, dass manche Maßnahmen unterbleiben oder
aufgeschoben werden (Auf Grund des Selbstschutzes, \enquote{Aufklärung} eines bewusstlosen Patienten,
\ldots{}) oder angepasst werden müssen.

{\scriptsize (Die Reihenfolge ist der jeweiligen Situation anzupassen!)}
\begin{itemize}
  \item Einschätzungsblock (\prettyref{m:einschaetzungsblock}?
  \item Beengende oder behindernde Kleidung entfernen bzw.
  öffnen
  \item Situationsgerechte Lagerung
  \item \E{Wärmeerhalt} oder \E{Kühlung}
  \item Angemessene Erhebung der Anamnese 
  \item Angemessene Untersuchung 
  \item Spezielle Maßnahmen gemäß Verdachtsdiagnose(n)
  \item Patientenidentifkation
  \item Dokumentation, Aufklärung
  \item Psychischer Beistand
  \item Verlaufskontrolle, Patientenbeobachtung, Monitoring
  \columnbreak
  \item Weiteres Vorgehen, je nach Bedarf und Situation:
  
  
 
  \begin{itemize}\label{tab:pam-weiterfuehrende}
    \item Transportentscheidung und ggfs. Transport an eine
    geeignete Einrichtung (Krankenanstalt).
    \item Notarzt-Nachforderung (bei Bedarf, z.\,B. Schmerztherapie,
    Aufklärung zwecks \E{Belassung}\index{Belassung} auf
    Patientenwunsch \E{trotz Behandlungsnotwendigkeit}, \ldots)
%     \item Unterlassung der Intervention wenn
%     \E{offensichtlich} kein sanitätsdienstliches Problem
%     besteht (Fehlalarmierung) 
%     Behandlungsbedarf/-wunsch
%     besteht und eine gesundheitliche Gefährdung des Patienten ausgeschlossen werden kann.
    \Ca{Rechtliche Hinweise beachten!}
    \item Patient begibt sich selbstständig in weitere
    Behandlung. Z.\,B. Bagatellverletzungen (z.\,B.
    oberflächliche Schürfwunde) ist es zulässig, dass sich
    der Patient selbständig in ärztliche Behandlung begibt.
    Das Anraten, eine ärztliche Behandlung aufzusuchen, ist
    zu dokumentieren und ggfs. vom Patienten per
    Unterschrift bestätigen zu lassen (Revers). Über Risiken
    muss aufgeklärt (und diese Aufklärung ebenfalls
    dokumentiert) werden.
  \end{itemize}
  \item Ggfs. Übergabe an weiterbehandelndes Personal
\end{itemize}

% \Commentary[Immer durchzuführende Maßnahmen, Belassung]
%   {Die Frage, wie nichtärztliches Personal
%   bei einer Transportverweigerung oder vermeintlich nicht
%   vorliegender Behandlungsbedürftigkeit/Spitalspflichtigkeit
%   vorzugehen hat, ist nach wie vor nicht
%   zufriedenstellend beantwortet.\ParBugfix
%   Zur rechtlichen Lage in Österreich: Die Untersuchung auf
%   das Vorliegen oder Nichtvorliegen von
%   körperlichen und psychischen Krankheiten oder Störungen
%   ist im \protect\S\,2 Abs.~2 Z.~1 ÄrzteG geregelt und hat keine
%   Entsprechung im SanG; somit fällt die Feststellung, ob
%   eine Behandlungsbedürftigkeit besteht unter den
%   Ärztevorbehalt, d.\,h. ist den Ärzten vorbehalten. \ParBugfix
%   \protect\Speech{\enquote{Ob Behandlungsbedarf vorliegt, kann nur
%   von einem Arzt entschieden werden. Denn die dafür erforderliche
%  Diagnose von (behaupteten) Krankheitszuständen fällt unter
%  \protect\S\,2 Abs.~2 Z.~1  ÄrzteG und ist daher den Ärzten
%  vorbehalten.}} OGH (4 Ob 36/10p)
%  \ParBugfix
%  Das \enquote{selbstständige Begeben in ärztliche Behandlung}
%  stellt keine Verneinung des Vorliegens einer Erkankung
%  oder gesundheitlichen Störung sowie deren 
%  grundsätzlicher Behandlungsnotwendigkeit dar.} % Commentary
 }

\pagestyle{empty}\loadgeometry{WholePage}\clearpage
\noindent\begin{minipage}{\linewidth}
% \noindent{\sffamily{\color{ubuntured}%
% {\bfseries\Huge\color{DefaultA}
% Patientenmanagement}%\bigskip 
% 
% Arbeits- und Ausbildungsstandards} {\color{darkBlue}für den Sanitätsdienst}} | {\AASS} 
% \hfill\textsf{\DocVersion}

% \bigskip

\Layout{57}{}{}{}
{\includegraphics[width=\linewidth]{./images/emhofer-josef-aass/PAMneuAugust2012}}
{Graphische Übersicht: \EEE{Standardisierte Patientenversorgung}}
{}

\vfill

\Layout{57}{}{}{}{%
\small
\begin{tabularx}{\linewidth}{XXX}
\addlinespace\toprule\addlinespace
\noindent\TabHeading{Massive Störung}
&
\noindent\TabHeading{Alarm-Symptome}
&
\noindent\TabHeading{Alarm-Diagnosen}
\\\addlinespace\midrule\addlinespace
Bewusstsein

Atemweg / Atmung 

Kreislauf 

(BeABC)
&
Atemnot, die sich nicht bessert 

Brodelndes Atemgeräusch 

Thoraxschmerz  

Schocksymptome 

entgleiste Vitalwerte

schwere Verletzungen

Hirndruckzeichen

{\dots} 
&
Herzinfarkt

Kardiales Lungenödem

Status Epilepticus

Beckenfraktur

Rückenmarksverletzung

\ldots
\\\addlinespace\bottomrule
\end{tabularx}
}{{\color{DefaultRed}\bfseries{\VarFlag}} \Ca{Alarmzeichen:}
Wann ist eine vitale Bedrohung
wahrscheinlich? Eine Übersicht.\label{tab:alarmzeichen}}{}

\vfill
\noindent{\sffamily\small{\textcopyright} \today\ \ARGEAASS \hfill \url{http://www.aass.at}}
\end{minipage}
\clearpage
% 
% \vspace{-\baselineskip}~
% ~
\vfill





% % M %%%%%%%%%%%%%%%%%%%%%%%%%%%%%%%%%%%%%%%%%%%%%%%%%%% M %
% % 
% \ProcedurePrintDefault{}
\MassnahmenWide{Standardmaßnahmen}{Einschätzungsblock}
{\label{m:ersteinschaetzung}\label{m:einschaetzungsblock}}
{%
\CorrectionItemizeBox\renewcommand{\itemhook}{%
%   \setlength{\topsep}{0pt}%
  \setlength{\itemsep}{0.004\baselineskip}
%   \setlength{\parskip}{\baselineskip}
  \setlength{\leftmargin}{1.4em}
  \setlength{\labelwidth}{0.9em}
}
\noindent\begin{minipage}[t]{\linewidth -
{2\fboxrule} - {2\fboxsep}}
\noindent\begin{minipage}[t]{\linewidth}
{\vspace{-2\baselineskip}%
\scriptsize
% \begin{tabulary}{\linewidth}[H]{cp{8em}p{8.5cm}}
\begin{tabularx}{\linewidth}[H]{p{1.5cm}XX}
% \toprule\addlinespace
&
\TabHeading{Beurteilung / Untersuchung}
&
\TabHeading{Typische Sofortmaßnahmen}
\\\addlinespace\toprule\addlinespace
\noindent\TabHeading{Szeneüberblick m. Schutz}
&
\begin{minipage}[t]{\linewidth}
\begin{itemize}
  \item Sicherheit, Selbstschutz; Patientenanzahl; 
  \item Umgebung, Ort,
Zeit, Gefahrenzonen, Umstände,
\item Weitere Resourcen erforderlich? Unfallmechanismus,
Großschaden
\item Lagemedlung erforderlich
\item \EE{Trauma?} \EE{Mechanismus?}
\end{itemize}
\end{minipage}
&
\noindent\color{Indigo}\begin{minipage}[t]{\linewidth}
\begin{itemize}
  \item GAS-Maßnahmen 
  \item Lagemeldung
  \item Schutzausrüstung
  \item Nachforderung weiterer Kräfte
\end{itemize}
\end{minipage}
\\\addlinespace\toprule\addlinespace
\noindent\TabHeading{Eindruck}
&
\begin{minipage}[t]{\linewidth}
\begin{itemize}
  \item Alter, Geschlecht, \EE{Hautfarbe}, Gesichtsausdruck,
  \EE{Haltung}, spontane Bewegungen,  Sprache
  \item Offensichtliche Verletzungen, Blutungen
  \item Sonstige \EE{Auffälligkeiten}
\end{itemize}
\end{minipage}
&
\noindent\color{Indigo}\begin{minipage}[t]{\linewidth}
\begin{itemize}
  \item Manuelle HWS-Fixierung
  \item Stillung von starken Blutungen
  \item Bewegungsverbot
\end{itemize}
\end{minipage}
\\\addlinespace\midrule\addlinespace
\noindent\TabHeading{Bewusstsein}
&
\begin{minipage}[t]{\linewidth}
\begin{itemize}
  \item \EE{Bewusstseinsgrad} (WASB)
\end{itemize}
\end{minipage}
&
\noindent\color{Indigo}\begin{minipage}[t]{\linewidth}
\begin{itemize}
  \item NA-Nachforderung
\end{itemize}
\end{minipage}
\\\addlinespace\midrule\addlinespace
\noindent\TabHeading{Hauptbeschwerde}
&
\begin{minipage}[t]{\linewidth}
\begin{itemize}
  \item \EE{Berufungsgrund} und
  \EE{Leitsymptome}
\end{itemize}
\end{minipage}
&
\noindent\color{Indigo}\begin{minipage}[t]{\linewidth}
\end{minipage}
\\\addlinespace\toprule\addlinespace
\noindent\T{A}

\TabHeading{Atemweg}
&
\begin{minipage}[t]{\linewidth}
\begin{itemize}
  \item *Inspektion der Atemwege (\EE{Mund}, \EE{Nase})
  \item \EE{Atemgeräusche}
\end{itemize}
\end{minipage}
&
\noindent\color{Indigo}\begin{minipage}[t]{\linewidth}
\begin{itemize}
  \item Absaugung, Fremdkörper entfernen
  \item Kopf überstrecken, Esmarch-Handgriff
  \item HWS-Immobilisation manuell oder mit Schiene
  \item \Z{Erweiterte Maßnahmen}
\end{itemize}
\end{minipage}
\\\addlinespace\midrule\addlinespace
\noindent\T{B}

\TabHeading{Atmung}
&
\begin{minipage}[t]{\linewidth}
\begin{itemize}
  \item Schätzen: \EE{Atemfrequenz}, \EE{-tiefe}
  \item *\EE{Sp\ce{O2}}
  \item Atemgeräusche
  \item Inspektion der \EE{Thorax-Atembewegungen}, ggfs. 
  Atemprobe
  \item \E{Zeichen der
  Atemnot}: Hautfarbe, Bauch-/Thorax\-bewegungen,
  Anstrengung beim Atmen, Atemhilfsmuskulatur, \ldots
  \item \Z{(Auskultation: Seitenvergleich, Lungenbasen)}
\end{itemize}
\end{minipage}
&
\noindent\color{Indigo}\begin{minipage}[t]{\linewidth}
\begin{itemize}
  \item Reanimation
  \item Assistierte Beatmung (bei AF \textless\,8 oder
  \textgreater\,30\PerMin*)
  \item \ce{O2}-Gabe
  \item Situationsgerechte Lagerung
  \item \Z{Erweiterte Maßnahmen\\(Entlastung
  Spannungspneumothorax, \ldots )}
\end{itemize}
\end{minipage}
\\\addlinespace\midrule\addlinespace
\noindent\T{C}

\TabHeading{Kreislauf}
&
\begin{minipage}[t]{\linewidth}
\begin{itemize}
  \item \EE{Hautfarbe}, \EE{-temperatur}
  \item Schweiß
  \item \EE{Rekap}-Zeit
  \item \EE{Radialispuls}: Stärke, Rhythmus, Abschätzen der
  Frequenz
\end{itemize}
\end{minipage}
&
\noindent\color{Indigo}\itshape\begin{minipage}[t]{\linewidth}
\begin{itemize}
  \item Blutstillung
  \item Situationsgerechte Lagerung
\end{itemize}
\end{minipage}
\\\addlinespace\cmidrule{2-3}\addlinespace
\noindent\TabHeading{*STU}
&
\begin{minipage}[t]{\linewidth}
\begin{itemize}
  \item Inspektion
  und Abtasten von
  \begin{inparaenum}
    \item \EE{Kopf} inkl. Ohren,
    \item \EE{Hals},
    \item \EE{Thorax},
    \item \EE{Bauch},
    \item \EE{Becken} und
    \item \EE{Oberschenkel}
  \end{inparaenum}
\end{itemize}
\end{minipage}
&
\noindent\color{Indigo}\begin{minipage}[t]{\linewidth}
\begin{itemize}
  \item Blutstillung
  \item Situationsgerechte Lagerung
\end{itemize}
\end{minipage}
\\\addlinespace\midrule\addlinespace
\noindent\T{D}

\TabHeading{Schnelle\newline Untersuchung}
&
\begin{minipage}[t]{\linewidth}
\begin{itemize}
  \item Vitalwerte (\EE{HF}, \EE{RR})
  \item Neurologisch:
  \begin{inparaitem}
    \item \EE{Orientierung} oder *\EE{GCS}
    \item \EE{Pupillen}
    \item \EE{Kraftprobe OE}, *UE
    \item *Kann der
    Patient \EE{Hände} und \EE{Füße} spüren und aktiv
    bewegen?
  \end{inparaitem}
  \item *\EE{Blutzuckermessung}
\end{itemize}
\end{minipage}
&
\noindent\color{Indigo}\begin{minipage}[t]{\linewidth}
\begin{itemize}
  \item Transportentscheidung (eilig, Rendez-vous, \ldots)
\end{itemize}
\end{minipage}
\\\addlinespace\midrule\addlinespace
\noindent\T{E}

\TabHeading{Erweiterte\newline Untersuchung}
&
\begin{minipage}[t]{\linewidth}
\begin{itemize}
  \item (Fremd-) \EE{Anamnese} (SAMPLE, OPQRST)
  \item Einschätzen der \EE{Umgebung}
  \item Weitere relevante \EE{Untersuchungen}
  \item Arbeitsdiagnose erstellen und Differentialdiagnosen
  ausschließen
\end{itemize}
\end{minipage}
&
\noindent\color{Indigo}\begin{minipage}[t]{\linewidth}
\begin{itemize}
  \item Transportentscheidung 
  \item Spezielle Maßnahmen gemäß Verdachtsdiagnose
\end{itemize}
\end{minipage}
% \\\addlinespace\bottomrule\addlinespace
% \TabHeading{Beurteilung}
% &
% \begin{minipage}[t]{\linewidth}
% \begin{itemize}
%   \item \EE{Einschätzen der vitalen Bedrohung}
%   ständig während des gesamten Blocks
%   \item Vorläufige Verdachtsdiagnose formulieren
% \end{itemize}
% \end{minipage}
% &
% \noindent\color{Indigo}\begin{minipage}[t]{\linewidth}
% \begin{itemize}
%   \item {\TxMassVit}
%   \item Reihenfolge der weiteren Untersuchungen und
%   Maßnahmen festlegen
% \end{itemize}
% \end{minipage}
\\\addlinespace\bottomrule
\end{tabularx}

\vspace{-0.5\baselineskip}\begin{center}
  \EE{Der Einschätzungsblock muss in regelmäßigen Abständen
in angemessenem Umfang wiederholt
werden (Verlaufskontrolle)!}

{\tiny Mit einem * versehene Punkte werden nur durchgeführt,
wenn sie \E{situationsentsprechend} sind. \enquote{Typische
Sofortmaßnahmen} sind als \E{Beispiele} zu verstehen.
Bei \Ca{absolut zeitkritischen Patienten}, bei denen bereits die 
ABC-Einschätzung ergibt dass ein Transport nicht aufschiebbar ist, 
kann u.\,U. schon \Ca{nach C eine vorgezogene Strategie- 
und Transportentscheidung} getroffen werden. 
D und E sollen dann, sofern möglich, während des 
Transportes erfolgen. }
\end{center}
}
\end{minipage}
\end{minipage}\vspace{-\baselineskip}
}

\vfill
\restoregeometry\pagestyle{fancy}

\Layout{57}{}
{}
{}
{%
\begin{minipage}{\linewidth}
\FontTableBase\footnotesize
{\sffamily\FontTableBase\footnotesize\CorrectionItemizeBox\renewcommand{\itemhook}{%
%   \setlength{\topsep}{0pt}%
  \setlength{\itemsep}{0.005\baselineskip}
%   \setlength{\parskip}{\baselineskip}
  \setlength{\leftmargin}{1.4em}
  \setlength{\labelwidth}{0.9em}
}
\begin{tabulary}{\linewidth}[H]{lCL}
\toprule\addlinespace
{\multirow{6}{2cm}{\TabHeading{Anamnese}}}
 &
{\sffamily\textbf S}
 &
Symptome {\&} Schmerzen: OPQRST

{\tiny(Lokalisation, Stärke,  Qualität, Zeit/Verlauf, 
 Ausstrahlung, Linderung/Verstärkung)}
\\\addlinespace
 &
{\sffamily\textbf A}
 &
Allergien/Allergische Reaktion (auch
Allergien auf Medikamente)
\\\addlinespace
 &
 {\sffamily\textbf M}
 &
Medikamente {\tiny(Nimmt Patient Medikamente?
Wogegen sind diese?)}
\\\addlinespace
 &
{\sffamily\textbf P}
 &
Patientengeschichte: Krankheiten {\tiny(chronische, frühere
K.)}, Operationen, \ldots 
\\\addlinespace
 &
 {\sffamily\textbf L}
 &
Letzte ... (zB letzte Mahlzeit, letzte
Regelblutung, letzter Spitalsaufenthalt)
\\\addlinespace
 &
 {\sffamily\textbf E}
 &
Ereignisse vor Notfalleintritt {\tiny(zB
Anstrengung vor Brustschmerz, Sturz, ...)}
\\\addlinespace\midrule\addlinespace

{\multirow{8}{2cm}{\TabHeading{Diagnostik}}}
 &
{\sffamily AF}
 &
Atemfrequenz auszählen
\\\addlinespace
 &
({\sffamily Sp\ce{O2}})
 &
Sauerstoffsättigung im Blut, Monitoring
\\\addlinespace
 &
{\sffamily  BZ}
 &
Blutzucker {\tiny(Wann war die letzte
Mahlzeit/Insulinspritze? Diabetes bekannt?)}
\\\addlinespace
 &
{\sffamily Temp.}
 &
Körpertemperatur/Fieber
\\\addlinespace
 &
 {\sffamily\small Neuro\-status}
 &
(Bewusstseinsgrad), Orientierung,
Pupillen, Halbseitenzeichen, Herdblick, Meningismus,
retrograde Amnesie
\\\addlinespace
 &
{\sffamily\small Trauma\-check}
 &
Abnorme Beweglichkeit, DMS an allen
Extremitäten (Nagelbettprobe), Wunden, paradoxe Atmung,
Abwehrspannung des Bauches, ...
\\\addlinespace
 &
{\sffamily\small Körp. Unt.}
&
Sonstige körperliche Untersuchungen: Abdomen abtasten,
Körperstellen inspizieren, \ldots
\\\addlinespace
 &
{\sffamily\small \ldots}
&
EKG, Kapnometrie, \ldots
\\\addlinespace\midrule\addlinespace
\TabHeading{Maßnahmen}
&
&
\begin{minipage}{\linewidth}
\begin{itemize}
  \item (Lagerung)
  \item Verband, Blutstillung
  \item Schienung
  \item psychischer Beistand
  \item \EE{Wärmeerhaltung / Kühlung}
  \item Entscheidung Notarzt-Nachforderung wegen
  Schmerztherapie oder Aufklärung/Belassung
  \item {\TxMassOGabe}
  \item Regelmäßige \EE{Wiederholung des
  Einschätzungsblockes} {\tiny(in angemessenem Umfang)}
  \item \EE{Immer durchzuführende Standardmaßnahmen}
  (\prettyref{m:standardmassnahmen-immer})
  \item \ldots
\end{itemize}
\end{minipage}
\\\addlinespace\midrule\addlinespace
\TabHeading{Assistenz}
&
&
\begin{minipage}{\linewidth}
\begin{itemize}
  \item EKG anlegen
  \item Infusion/periphervenösen Zugang vorbereiten
  \item Intubation vorbereiten
  \item Medikamente vorbereiten
  \item \ldots
\end{itemize}
\end{minipage}
\\\addlinespace\midrule\addlinespace
\TabHeading{Transport}
&
&
\begin{minipage}{\linewidth}
\begin{itemize}
  \item Welches Spital/Bett buche ich ab? {\tiny(Arbeitsunfall?
  Welche Abteilung fahre ich an? Ist der Patient bereits in
  einem Spital in Behandlung? Spitalsbehandlung vor kurzer
  Zeit (\enquote{Reklamation})?)}
  
  Brauche ich eine Spezialabteilung? {\tiny(Schockraum, Stroke
  Unit, PTCA, Verbrennungsstation)}
  \item Übergabe {\tiny(Anamnese, Befunde, Maßnahmen, Hinweise)}
\end{itemize}
\end{minipage}

\\\addlinespace\bottomrule
\end{tabulary}}

{\scriptsize\EE{Fett} gedruckte Punkte haben eine besonders hohe Priorität, (eingeklammerte) Punkte sollten schon als Teil des Einschätzungsblockes abgehandelt worden sein.}
\end{minipage}
}
{Kurzübersicht: \EEE{Weiterführende Untersuchungen und Anamnese}.
Die Reihenfolge ist je nach Patient
unterschiedlich!\label{tab:checkliste-patientenmanagement}\label{tab:pam-weiterfuehrende}.}
{}






% % M %%%%%%%%%%%%%%%%%%%%%%%%%%%%%%%%%%%%%%%%%%%%%%%%%%% M %
% % 
% \ProcedurePrintDefault{}
\MassnahmenNG{Standardmaßnahmen}{Standardmaßnahmen bei vital bedrohten Patienten}{\label{m:standardmassnahmen-vital}}{
\label{topic:standardmassnahmen-vital}%

\begin{enumerate}
  \item \EE{Situationsgerechte Lagerung}
  \item \EE{Sauerstoffgabe}
  (\prettyref{m:sauerstoffberieselung}): je nach Indikation,
  allgemeiner Zielwert: Sp\ce{O2} von \VonBis{94}{98}\PC*
  \item Ggfs. \EE{Notarzt-Nachforderung} :
  mit kurzer Begründung
  
  \Ca{Bei manchen Notfällen keine Notarztnachforderung wenn
  der Patient zeitkritisch ist, siehe jeweilige spezielle
  Maßnahmen!}
  \item \EE{Engmaschiges, bestmögliches Monitoring}: Je nach
  vorhandenem Material
  
  RR, HF, Pulsoxymetrie, EKG,
  Sitzwache/Patientenbeobachtung, etc.\footnote{Zuerst werden die Patientenbeobachtung
  und -- wenn vorhanden -- die
  Pulsoxymetrie eingesetzt. Überwachungsgeräte, deren Anlage
  zeitintensiv ist (EKG, \ldots), sollen erst dann verwendet
  werden, wenn alle dringlicheren Maßnahmen durchgeführt und
  der Einschätzungsblock beendet ist! Die
  Patientenbeobachtung bleibt immer ein wesentlicher Teil
  des Monitorings!}
  \item \EE{Reanimationsbereitschaft} herstellen
  (\prettyref{m:reanimationsbereitschaft})
\end{enumerate}

{\scriptsize(Reihenfolge zählt!)}
}




% % M %%%%%%%%%%%%%%%%%%%%%%%%%%%%%%%%%%%%%%%%%%%%%%%%%%% M %
% % 
% \ProcedurePrintDefault{}
\MassnahmenNG{Standardmaßnahmen}
{Reanimationsbereitschaft}
{\label{m:reanimationsbereitschaft}}
{%
\begin{enumerate}
  \item Ggfs. Platz schaffen\footnote{z.\,B. im Wohnzimmer den Couchtisch, Hocker etc. zur Seite schieben}
  \item Geräte vorbereiten:
\begin{enumerate}
    \item \EE{Beatmungsbeutel}
    \item \EE{Absaugeinheit} und passenden Katheter  in
    Griffweite stellen. Steriles Material bleibt verpackt.
    \item \EE{Defibrillator} in Griffweite stellen,
    Elektroden und Verbrauchsmaterial bleiben verpackt.
  \end{enumerate}
\end{enumerate}

}





% % M %%%%%%%%%%%%%%%%%%%%%%%%%%%%%%%%%%%%%%%%%%%%%%%%%%% M %
% % 
% \ProcedurePrintDefault{}
\MassnahmenNG{Standardmaßnahmen}
{Einschätzung der Indikation zur Wirbelsäulenimmobilisation}
{\label{m:einschaetzung-indikation-ws-immobilisation}}
{%
% \marginnote{\includegraphics[width=\marginparwidth]{./images/pallinger-christoph-aass/Spiderstrap_32954-AASS-0112mm.jpg}}
Eine Immobilisation der Wirbelsäule ist notwendig bei
Traumapatienten, bei denen folgendes zutrifft:

\begin{enumerate}
  \item \EE{Typischer Unfallmechanismus}
  \begin{itemize}
    \item Hochgeschwindigkeitsunfall
    \item Sturz $\geq$ 3\Meter* oder 3fache Patientengröße
    \item Penetrierende Verletzungen im Bereich der
    Wirbelsäule
    \item Sportverletzungen im Kopf- oder Nackenbereich
    \item Trauma nach Sprung ins Wasser
    \item Stauchungsverletzung der HWS (z.\,B. Schlag auf
    Kopf)
    \item Suizidversuch durch Erhängen
    \item Bewusstloser Traumapatient
  \end{itemize}
  \item \EE{Unklarer Unfallmechanismus}, bei dem einer der
  folgenden Punkte zutrifft:
  \begin{itemize}
    \item Spinaler Schmerz oder Druckschmerz über der
    Wirbelsäule
    \item Auffälliger Befund bei Untersuchung von Motorik
    und Sensibilität\footnote{sofern nicht anders plausibel
    erklärbar}
    \item Patientenangaben kann nicht vertraut
    werden
    \begin{itemize}
      \item Stresssituation
      \item Schädel- oder zerebrale Verletzung
      \item Auffälliger Bewusstseinszustand
      \item Intoxikation
      \item Ablenkende Verletzungen
    \end{itemize}
  \end{itemize}
  \item \EE{Verdacht}
\end{enumerate}

Solange es noch unklar ist, ob eine WS-Immobilisation
notwendig ist, muss eine manuelle Fixierung der HWS
durchgeführt werden.

\cite{ITLSde6,ScholzNotfallmedizin2,SiewertChirurgie8}
}
\Commentary[Konsensus]{Review 2011-01-11}

\chapter{Medizinische Maßnahmen}




% % M %%%%%%%%%%%%%%%%%%%%%%%%%%%%%%%%%%%%%%%%%%%%%%%%%%% M %
% % 
% \ProcedurePrintDefault{}
\MassnahmenNG{Spezielle Maßnahmen}{Bewusstseinseintrübung}
{\label{m:bewusstseinstruebung}}{ 
% \paragraph{Maßnahmen
% bei getrübten Patienten}~

\begin{itemize}
  \item Lagerung: situationsgerecht je nach
  Verdachtsdiagnose, im Zweifel \EE{stabile
  Seitenlage}, bei hochschwangeren Frauen auf die
  \E{linke} Seite.
  \item \TxBeiVit \TxMassVitMKBes
  
  \begin{itemize}
    \item \EE{Monitoring}: Auf \EE{Bewußtseinsgrad}
    besonders achten!
  \end{itemize}
  
  
  \item Diagnostische Schwerpunkte:
  \begin{itemize}
    \item \EE{Ursachenforschung} -- Warum ist der
    Patient eingetrübt?
    \item (Fremd-)Anamnese, Umstände, Szenerie
    \item \EE{Neurocheck}, inkl. \EE{BZ}-Messung!
    \item \EE{Traumacheck}. Suche nach Primär- und
    Folgeverletzungen!
  \end{itemize}
\end{itemize}

}




% % M %%%%%%%%%%%%%%%%%%%%%%%%%%%%%%%%%%%%%%%%%%%%%%%%%%% M %
% % 
% \ProcedurePrintDefault{}
\MassnahmenNG{Spezielle Maßnahmen}{Bewusstlose und
soporöse Patienten}{\label{m:bewusstlosigkeit}}
{%
\begin{itemize}
  \item \TxMassVitMKBes
  \begin{itemize}
    \item Lagerung: stabile Seitenlage, bei hochschwangeren Frauen auf die
    \E{linke} Seite.
    \item \EE{Monitoring}: Auf \EE{Bewußtseinsgrad}
    besonders achten!
  \end{itemize}
  \item Diagnostische Schwerpunkte:
  \begin{itemize}
    \item \EE{Ursachenforschung} -- Warum ist der Patient bewusstlos oder
    soporös?
    \item (Fremd-)Anamnese, Umstände, Szenerie
    \item \EE{Neurocheck}, inkl. \EE{BZ}-Messung!
    \item \EE{Traumacheck}. Suche nach Primär- und
    Folgeverletzungen!
  \end{itemize}
\end{itemize}
}




% % M %%%%%%%%%%%%%%%%%%%%%%%%%%%%%%%%%%%%%%%%%%%%%%%%%%% M %
% % 
% \ProcedurePrintDefault{}
\MassnahmenNG{Spezielle
Maßnahmen}{Insuffiziente
Atmung (AF \textless\,8 oder
\textgreater\,30\PerMin*,
bzw. AZV zu niedrig)}{\label{m:atmung-insuffizient}}
{\Ca{Vitale Bedrohung}
\begin{itemize}
  \item \TxMassVitMK
  \item Assistierte Beatmung
  \item Ursachenforschung
\end{itemize}
}




% % M %%%%%%%%%%%%%%%%%%%%%%%%%%%%%%%%%%%%%%%%%%%%%%%%%%% M %
% % 
% \ProcedurePrintDefault{}
\MassnahmenNG{Spezielle
Maßnahmen}{Atemstillstand}{\label{m:atemstillstand}}{
\Cave{Achtung: Bei einem mit Atemstillstand vorgefundenen
Patienten ist grundsätzlich von einer Reanimation
auszugehen.}

\Z{Nur in sehr seltenen Fällen kann primär ein isolierter
Atemstillstand ohne gleichzeitigem Kreislaufstillstand
beobachtet werden. Auf diese Spezialfälle (z.\,B.
Opiatintoxikation) wird hier nicht weiter eingegangen. Auch
der iatrogene Atemstillstand (z.\,B. im Rahmen einer Narkose) wird an dieser Stelle nicht
behandelt.}}




\ProcedurePrintDefault{MR57--0B}


















% % M %%%%%%%%%%%%%%%%%%%%%%%%%%%%%%%%%%%%%%%%%%%%%%%%%%% M %
% % 
% \ProcedurePrintDefault{}
\MassnahmenWide{Spezielle Maßnahmen}{Unterlassung der
Reanimation} {\label{m:reanimation-unterlassung}}
{%
\Ca{Unter den folgenden Voraussetzungen darf (oder muss) eine
Reanimation unterlassen
werden:\index{Reanimation!Unterlassung
der}\index{Wiederbelebung|see{Reanimation}}}           
\begin{itemize}
  \item \EE{Verwesung}
  \item \EE{Fäulnis}
  \item \EE{Mumifizierung}
  \item \EE{Skelettierung}
  \item \EE{Verletzungen, die nicht mit dem Leben
  vereinbar sind}
  \item \EE{Tierfraß}\footnote{Unter
  \E{\enquote{Tierfraß}} versteht man den Verzehr eines Kadavers
  durch Kadaverfresser. Ein frischer Hundebiss oder
  Ähnliches fällt \E{nicht} darunter. Maden können
  sich im Gewebe von Lebenden einnisten und sind daher
  auch kein sicheres Todeszeichen.}
  \item (Verbindliche) \EE{Patientenverfügung}\footnote{\E{Patientenverfügung}: Bei Notfallversorgungen darf \EE{nicht nach einer Verfügung gesucht
  werden, wenn dadurch das Leben oder die Gesundheit des Patienten ernstlich
  gefährdet werden würde}.
  \EE{Aber:} Wenn \E{zweifelsfrei} eine verbindliche
  Patientenverfügung vorliegt (z.\,B. weil eine betreuende
  Pflegekraft diese in der Zwischenzeit gefunden und die
  Verbindlichkeit bestätigt hat), muss dieser Folge
  geleistet werden.
Liegt eine \E{beachtliche Patientenverfügung} vor, so ist im
\E{Einzelfall} zu entscheiden: Das Wohl des Patienten bleibt
oberstes Gebot und die Verfügung soll in die
Entscheidung einfließen.}\index{Patientenverfügung!Unterlassung der
  Reanimation} (\FullPrettyRef{AASS-topic:patientenverfuegung}; 
  Bis zur zweifelsfreien Klärung der Situation muss
  jedenfalls eine Reanimation begonnen bzw. fortgeführt
  werden!)
  \item (Abbruch der Reanimationsmaßnahmen aufgrund
  von \EE{Erschöpfung})\footnote{Erschöpfung darf im
  professionellen Umfeld keine Rolle spielen, für eine
  entsprechende Ablösung ist bei Bedarf zu sorgen.}
  \item Unvereinbarkeit mit
  \EE{Selbstschutz}\footnote{Es müssen jedoch alle Maßnahmen
  getroffen werden, die ohne Gefährdung des Selbstschutzes
  möglich sind. So wird zum Beispiel für Ersthelfer
  empfohlen, wenn
  eine Mund-zu-Mund-Beatmung nicht zumutbar ist, zumindest 
  eine Herzdruckmassage durchzuführen.}
\end{itemize}

}










% M %%%%%%%%%%%%%%%%%%%%%%%%%%%%%%%%%%%%%%%%%%%%%%%%%%% M %
% Herzinsuffizienz, symptomatisch
\ProcedurePrintDefault{MI50090C}


% M %%%%%%%%%%%%%%%%%%%%%%%%%%%%%%%%%%%%%%%%%%%%%%%%%%% M %
% Thoraxschmerzen, jetzt beschwerdefrei
\ProcedurePrintDefault{MI20091C}


% M %%%%%%%%%%%%%%%%%%%%%%%%%%%%%%%%%%%%%%%%%%%%%%%%%%% M %
% Akutes Koronarsyndrom
\ProcedurePrintDefault{MI24091C}


% M %%%%%%%%%%%%%%%%%%%%%%%%%%%%%%%%%%%%%%%%%%%%%%%%%%% M %
% 
\ProcedurePrintDefault{MR47091C}


% M %%%%%%%%%%%%%%%%%%%%%%%%%%%%%%%%%%%%%%%%%%%%%%%%%%% M %
% Kollaps/Synkope
\ProcedurePrintDefault{MR55--0C}


% M %%%%%%%%%%%%%%%%%%%%%%%%%%%%%%%%%%%%%%%%%%%%%%%%%%% M %
% Hypertensive Krise
\ProcedurePrintDefault{MI10911C}


% M %%%%%%%%%%%%%%%%%%%%%%%%%%%%%%%%%%%%%%%%%%%%%%%%%%% M %
% Hypertensiver Notfall
\ProcedurePrintDefault{MI10912C}


% M %%%%%%%%%%%%%%%%%%%%%%%%%%%%%%%%%%%%%%%%%%%%%%%%%%% M %
% Arterieller Gefäßverschluss einer Extremität
\ProcedurePrintDefault{MI74041C}


% M %%%%%%%%%%%%%%%%%%%%%%%%%%%%%%%%%%%%%%%%%%%%%%%%%%% M %
% Venöser Gefäßverschluss einer Extremität
\ProcedurePrintDefault{MI74042C}


% M %%%%%%%%%%%%%%%%%%%%%%%%%%%%%%%%%%%%%%%%%%%%%%%%%%% M %
% Mechanische Atemwegsverlegung
\ProcedurePrintDefault{MT17091C}


% M %%%%%%%%%%%%%%%%%%%%%%%%%%%%%%%%%%%%%%%%%%%%%%%%%%% M %
% Akuter Asthmaanfall
\ProcedurePrintDefault{MJ46--1C}


% M %%%%%%%%%%%%%%%%%%%%%%%%%%%%%%%%%%%%%%%%%%%%%%%%%%% M %
% COPD-Exazerbation
\ProcedurePrintDefault{MJ44190C}


% M %%%%%%%%%%%%%%%%%%%%%%%%%%%%%%%%%%%%%%%%%%%%%%%%%%% M %
% Lungenembolie
\ProcedurePrintDefault{MI26--0C}
% \MassnahmenNG{Spezielle
% Maßnahmen}{Lungenembolie}{\label{m:lungenembolie}}{
% \EE{Taktik}:
% \Speech{\Ca{Vitale Bedrohung!} Atemnot lindern, rasche medikamentöse Therapie, ärztliche Therapie sinnvoll}
% 
% \begin{itemize}
%   \item {\TxMassVitMKBes}
%   \begin{itemize}
%     \item Lagerung: Oberkörper hoch
%   \end{itemize}
%   \item Striktes Bewegungsverbot!
%   \item Beengende Kleidung öffnen
% \end{itemize}
% }





% M %%%%%%%%%%%%%%%%%%%%%%%%%%%%%%%%%%%%%%%%%%%%%%%%%%% M %
% Akutes Lungenödem, leicht
\ProcedurePrintDefault{MJ81--1C}


% M %%%%%%%%%%%%%%%%%%%%%%%%%%%%%%%%%%%%%%%%%%%%%%%%%%% M %
% Akutes Lungenödem, schwer
\ProcedurePrintDefault{MJ81--2C}


% M %%%%%%%%%%%%%%%%%%%%%%%%%%%%%%%%%%%%%%%%%%%%%%%%%%% M %
%
\ProcedurePrintDefault{MF45331C}


% M %%%%%%%%%%%%%%%%%%%%%%%%%%%%%%%%%%%%%%%%%%%%%%%%%%% M %
% Insult
\ProcedurePrintDefault{NI64--0C}



% M %%%%%%%%%%%%%%%%%%%%%%%%%%%%%%%%%%%%%%%%%%%%%%%%%%% M %
% Krampfender Patient
\ProcedurePrintDefault{NR56081C}





% M %%%%%%%%%%%%%%%%%%%%%%%%%%%%%%%%%%%%%%%%%%%%%%%%%%% M %
% Nicht-(mehr) krampfender Patient
\ProcedurePrintDefault{NR56082C}



% M %%%%%%%%%%%%%%%%%%%%%%%%%%%%%%%%%%%%%%%%%%%%%%%%%%% M %
% Hypoglykämie
\ProcedurePrintDefault{ME16020C}




% M %%%%%%%%%%%%%%%%%%%%%%%%%%%%%%%%%%%%%%%%%%%%%%%%%%% M %
% Hyperglykämisches Koma
\ProcedurePrintDefault{ME14012C}




% % M %%%%%%%%%%%%%%%%%%%%%%%%%%%%%%%%%%%%%%%%%%%%%%%%%%% M %
% % 
% \ProcedurePrintDefault{}
\MassnahmenNG{Allgemeine Maßnahmen}{Abdominalerkrankung}
{\label{m:am-abdominalerkrankungen}}
{
~

\begin{itemize}
  \item \EE{Vitale Bedrohung einschätzen.} \TxBeiVit
  \TxMassVitMK
  \item Lagerung: grundsätzlich \EE{bauchdeckenentspannend}
  (Knierolle, Beine angezogen), oder nach Wunsch des Patienten
  \item \EE{Nüchtern lassen} (in Hinblick auf möglicherweise
  bevorstehende Operation, nichts essen oder trinken lassen)
  %     \item Monitoring. Der Zustand kann sich abrupt
  %     verschlechtern.
  \item \EE{Transportentscheidung}: Situationsgerecht,
  z.\,B. Abt. f. Chirurgie oder Innere Medizin.
  
  Bei Kindern Abt. f. Kinderchirurgie oder Abt. f.
  Kinderheilkunde (je nach Verdachtsdiagnose und regionaler
  Regelung).
  
  Bei weiblichen Patienten Abt. für Gynäkologie (evtl. im
  \enquote{Hintergrund}) erwägen.
\end{itemize}
}








% M %%%%%%%%%%%%%%%%%%%%%%%%%%%%%%%%%%%%%%%%%%%%%%%%%%% M %
% Akutes Abdomen
\ProcedurePrintDefault{CR10000C}







% % M %%%%%%%%%%%%%%%%%%%%%%%%%%%%%%%%%%%%%%%%%%%%%%%%%%% M %
% % 
% \ProcedurePrintDefault{}
\MassnahmenNG{Spezielle
Maßnahmen}{Darmverschluß}{\label{m:darmverschluss}}{
\EE{Taktik}: \Speech{Unmittelbare Bedrohung je nach Vitalwerten. Nüchtern lassen, Differentialdiagnosen
abklären, Transport an geeignete Einrichtung}


\begin{itemize}
  \item Allgemeine Maßnahmen bei Abdominalerkrankungen
  (\prettyref{m:am-abdominalerkrankungen})
  \item \EE{Vitale Bedrohung einschätzen}, ggfs.
  \TxMassVitMKBes
  \begin{itemize}
    \item Lagerung: bauchdeckenentspannend, bzw. je nach
    Patientenwunsch
  \end{itemize}
  \item Patient nüchtern lassen.
  \item Transportentscheidung: Abt. f. Chirurgie
\end{itemize}
}






% % M %%%%%%%%%%%%%%%%%%%%%%%%%%%%%%%%%%%%%%%%%%%%%%%%%%% M %
% % 
% \ProcedurePrintDefault{}
\MassnahmenNG{Spezielle
Maßnahmen}{Appendizitis}{\label{m:appendizitis}}{
\EE{Taktik}:
\Speech{Differntialdiagnosen abklären, Transport an geeignete Abteilung}


\begin{itemize}
  \item Allgemeine Maßnahmen bei Abdominalerkrankungen
  (\prettyref{m:am-abdominalerkrankungen})
  \item Patient nüchtern lassen.
  \item Transportentscheidung: Abt. f. Chirurgie, Sonderfälle
  (Kinder, Frauen!) beachten!
\end{itemize}
}





% % M %%%%%%%%%%%%%%%%%%%%%%%%%%%%%%%%%%%%%%%%%%%%%%%%%%% M %
% % 
% \ProcedurePrintDefault{}
\MassnahmenNG{Spezielle
Maßnahmen}{Gallenkolik}{\label{m:gallenkolik}}{
\EE{Taktik}: \Speech{Rasche Schmerztherapie (veranlassen)}


\begin{itemize}
  \item Allgemeine Maßnahmen bei Abdominalerkrankungen
  (\prettyref{m:am-abdominalerkrankungen})
  \item Rasche Schmerztherapie veranlassen: Je nach
  Situation, Notarzt oder rascher Transport in eine
  geeignete Einrichtung
  \item Patient nüchtern lassen
  \item Transportentscheidung: Abt. f. Chirurgie
\end{itemize}
}




% % M %%%%%%%%%%%%%%%%%%%%%%%%%%%%%%%%%%%%%%%%%%%%%%%%%%% M %
% % 
% \ProcedurePrintDefault{}
\MassnahmenNG{Spezielle
Maßnahmen}{Blutungen
des Verdauungstrakts}{\label{m:blutung-verdauungstrakt}}{
\EE{Taktik}: \Speech{\Ca{Vitale Bedrohung bei hohem Blutverlust!}
Schockbekämpfung, zügiger Transport}


\begin{itemize}
  \item \TxMassVitMK
    \item Lagerung:
    Situationsgerecht, je nach Blutungsquelle und
    Patientenzustand
  \item Allgemeine Maßnahmen der Schockbekämpfung
  (\prettyref{m:am-schock})
  \item Ggfs. spezielle Maßnahmen bei hypovolämischem Schock
  (\prettyref{AASS-topic:schock-hypovolaemischer})
  \item Transportentscheidung: Abteilung für Chirurgie,
  Voranmeldung!
\end{itemize}
}





% % M %%%%%%%%%%%%%%%%%%%%%%%%%%%%%%%%%%%%%%%%%%%%%%%%%%% M %
% % 
% \ProcedurePrintDefault{}
\MassnahmenNG{Spezielle
Maßnahmen}{Gastroenteritis}{\label{m:gastroenteritis}}{
\begin{itemize}
  \item Allgemeine Maßnahmen bei Abdominalerkrankungen
  (\prettyref{m:am-abdominalerkrankungen})
  \item Patient nüchtern lassen
  \item Transportentscheidung: Abt. f. Innere Medizin; bei
  Kindern Abt. f. Kinderheilkunde
  \item Strikte Desinfektionsmaßnahmen!
\end{itemize}
}




% % M %%%%%%%%%%%%%%%%%%%%%%%%%%%%%%%%%%%%%%%%%%%%%%%%%%% M %
% % 
% \ProcedurePrintDefault{}
\MassnahmenNG{Spezielle Maßnahmen}{Gastritis,
Duodenitis}{\label{m:gastro-duodenitis}}{
\begin{itemize}
  \item Allgemeine Maßnahmen bei Abdominalerkrankungen
  (\prettyref{m:am-abdominalerkrankungen})
  \item \EE{Transportentscheidung}: 
  Abt. f. Chirurgie
%   
%   Bei Blutungszeichen (Teerstuhl, kaffeesatzartiges
%   Erbrechen): Abt. f. Chirurgie
\end{itemize}
}






% % M %%%%%%%%%%%%%%%%%%%%%%%%%%%%%%%%%%%%%%%%%%%%%%%%%%% M %
% % 
% \ProcedurePrintDefault{}
\MassnahmenNG{Spezielle Maßnahmen}{Akute
Pankreatitis}{\label{m:pankreatitis-akut}}{
\begin{itemize}
  \item Allgemeine Maßnahmen bei Abdominalerkrankungen
  (\prettyref{m:am-abdominalerkrankungen})
  \item \EE{Transportentscheidung}: Abt. f. Chirurgie
\end{itemize}
}





% % M %%%%%%%%%%%%%%%%%%%%%%%%%%%%%%%%%%%%%%%%%%%%%%%%%%% M %
% % 
% \ProcedurePrintDefault{}
\MassnahmenNG{Spezielle
Maßnahmen}{Mesenterialinfarkt}{\label{m:mesenterialinfarkt}}{
\begin{itemize}
  \item Allgemeine Maßnahmen bei Abdominalerkrankungen
  (\prettyref{m:am-abdominalerkrankungen})
  \item Nüchtern lassen!
  \item \EE{Transportentscheidung}: Abt. f.
  Chirurgie
\end{itemize}
}





% % M %%%%%%%%%%%%%%%%%%%%%%%%%%%%%%%%%%%%%%%%%%%%%%%%%%% M %
% % 
% \ProcedurePrintDefault{}
\MassnahmenNG{Spezielle
Maßnahmen}{Bauchfellentzündung}{\label{m:bauchfellentzuendung}}{
\begin{itemize}
  \item Maßnahmen wie bei akutem Abdomen
  (\FullPrettyRef{m:akutes-abdomen}), und sonst abhängig von
  der Ursache
\end{itemize}
}





% % M %%%%%%%%%%%%%%%%%%%%%%%%%%%%%%%%%%%%%%%%%%%%%%%%%%% M %
% % 
% \ProcedurePrintDefault{}
\MassnahmenNG{Spezielle
Maßnahmen}{Nierenkolik}{\label{m:nierenkolik}}{ Im
Vordergrund bei der Versorgung steht der rasche zügige Transport an eine Abteilung für Urologie. Wenn die Schmerzen
sehr massiv sind, muss eventuell eine Schmerztherapie durch
den Notarzt erfolgen.
\begin{itemize}
  \item Lagerung nach Wunsch des Patienten
  \item Bei unerträglichen Schmerzen Notarzt zur
  Schmerztherapie beiziehen
  \item Transportentscheidung: Abt. für Urologie
\end{itemize}
}





% % M %%%%%%%%%%%%%%%%%%%%%%%%%%%%%%%%%%%%%%%%%%%%%%%%%%% M %
% % 
% \ProcedurePrintDefault{}
\MassnahmenNG{Spezielle Maßnahmen}{Akutes
Harnverhalten}{\label{m:harnverhalten}}
{
\begin{itemize}
    \item Zügiger Transport in eine Urologische Ambulanz
    \item \Z{vorsichtig ablassen über Harnkatheter oder
    Blasenpunktion (Spital)}
\end{itemize}
}






% % M %%%%%%%%%%%%%%%%%%%%%%%%%%%%%%%%%%%%%%%%%%%%%%%%%%% M %
% % 
% \ProcedurePrintDefault{}
\MassnahmenNG{Spezielle
Maßnahmen}{MRSA-Transport}{\label{m:mrsa-transport}}{

\label{bkm:RefHeading40262819}

\begin{itemize}
  \item Informationen über MRSA-Trägerschaft sowie Art/Ort der MRSA-Besiedelung
  \item Patientenvorbereitung und Transport
  \begin{itemize}
    \item Wunden müssen verbunden und gut abgedeckt sein
    \item Bei Besiedelung der Atemwege trägt der Patient einen Mund-/ Nasenschutz
    \item Vor dem Transport führt der Patient eine hygienische
    Händedesinfektion durch
  \end{itemize}
  \item Allgemeine Hygienemaßnahmen
  \begin{itemize}
    \item Einweghandschuhe und Schutzkittel
    \item Nach Ablegen sofortige hygienische Händedesinfektion
    \begin{itemize}
      \item Nach Abschluss des Patiententransportes sind alle Materialien,
      Geräte, Instrumente und Flächen, welche direkten Kontakt mit dem
      Patienten hatten, gemäß Hygieneplan zu
      desinfizieren.
      \item Alle waagerechten Oberflächen des Fahrzeuginnenraumes sind mit
      einem Flächendesinfektionsmittel in Form einer
      Scheuerwischdesinfektion desinfizierend zu reinigen.
      \item Abschließend führt das Personal eine hygienische
      Händedesinfektion durch.
      \item An Keimvehikel denken (Kugelschreiber, Haltegriffe, Stethoskop \ldots )!
    \end{itemize}
  \end{itemize}
  \item Nach Abschluss der vorangegangenen Hygiene- und
  Desinfektionsmaßnahmen ist das Fahrzeug wieder einsatzbereit.
\end{itemize}

Für die anderen Problemkeime \Z{(VRSA, VRE, ESBL,
{\dots})} gelten sinngemäß die gleichen
Verfahrensweisen.
}





% % M %%%%%%%%%%%%%%%%%%%%%%%%%%%%%%%%%%%%%%%%%%%%%%%%%%% M %
% % 
% \ProcedurePrintDefault{}
\MassnahmenNG{Spezielle Maßnahmen}{Nadelstichverletzung}
{\label{m:nadelstichverletzung}}
{
\Z{Basiernd auf \cite{OeAG2007}.}
\begin{itemize}
    \item Erstmaßnahmen durchführen: \nocite{OeAG2007}
    \begin{itemize}
        \item bei Stich- bzw.Schnittverletzungen:
        Blutung durch Kompression des umgebenden
        Gewebes induzieren. Danach desinfizieren und einen
        in Desinfektionsmittel getränkten Tupfer 10\Min*
        auf der Wunde belassen.
        \item bei Spritzern mit Blut oder Sekreten in Auge oder Mund:
        Schleimhaut mit Wasser u./o. Schleimhautdesinfektionsmittel spülen
    \end{itemize}
    \item Leitstelle informieren
    \item ${\rightarrow}$ \EE{Umgehend in ein
    geeignetes Krankenhaus} (Von der Leitstelle erfragen).
    \item Patient: Zielspital informieren:
    \begin{itemize}
        \item Blut sichern für \Abk{HIV}- und Hepatitis-Serologie
    \end{itemize}
    \item  \Ca{Unfallmeldung}  (wichtig!), Blutabnahme
    \item ggf. \Abk{HIV}-Postexpositionsprohylaxe (\Abk{PEP})
    \begin{itemize}
        \item kurzfristige Therapie (4 Wochen)
        \item Einzelfall muss immer individuell beurteilt werden
        (Nutzen/Risiko), starke Nebenwirkungen.
        \item Muss innerhalb von 2\Hour* begonnen werden ${\rightarrow}$ daher
        \underline{\EE{sofort}} in ein geeignetes Krankenhaus!
    \end{itemize}
    \item Meldung an Betriebsleitung
    \item \EE{Unfallmeldung (wichtig!) an die \Abk{AUVA}}
\end{itemize}

\Cave{Unfallmeldung nicht vergessen!

Sofort in ein geeignetes Krankenhaus! (in Wien AKH od. OWS)}
}












% % M %%%%%%%%%%%%%%%%%%%%%%%%%%%%%%%%%%%%%%%%%%%%%%%%%%% M %
% % 
% \ProcedurePrintDefault{}
\MassnahmenNG{Spezielle
Maßnahmen}{Meningitis}{\label{m:meningitis}}{
\begin{itemize}
  \item Schutzmaske für Personal und Patient (wenn zumutbar)
  \item Rückmeldung an Leitstelle vor Transport, Voranmeldung im Krankenhaus
  \item Patient erst in das Spitalsgebäude bringen, wenn
  Ziel und Weg dorthin klar ist
  \item Desinfektionsmaßnhamen gem. Hygieneplan
  \item Personal bekommt im Spital evtl. eine Antibiotika-Prophylaxe
  \item Wechsel der Dienstkleidung
\end{itemize}

}







% % M %%%%%%%%%%%%%%%%%%%%%%%%%%%%%%%%%%%%%%%%%%%%%%%%%%% M %
% % 
% \ProcedurePrintDefault{}
\MassnahmenNG{Spezielle
Maßnahmen}{Tuberkulose}{\label{m:tuberkulose}}{
\begin{itemize}
  \item Bei Verdacht auf eine offene Tuberkulose
  Schutzmasken verwenden: FFP~3 für das Personal, \enquote{OP-Maske} für den
  Patienten, wenn zumutbar
  \item \EE{Meldung an Leitstelle \& Betriebsleitung!}
  \item Desinfektion laut Hygieneplan
\end{itemize}
}






% % M %%%%%%%%%%%%%%%%%%%%%%%%%%%%%%%%%%%%%%%%%%%%%%%%%%% M %
% % 
% \ProcedurePrintDefault{}
\MassnahmenNG{Spezielle
Maßnahmen}{Pneumonie}{\label{m:pneumonie}}{

\begin{itemize}
  \item Symptomatisch
  \item Wärmeerhalt
  \item Bei Fieber Meningismus-Zeichen prüfen!
  \item Transportentscheidung: Abt. f. Innere Medizin
\end{itemize}
}





% % M %%%%%%%%%%%%%%%%%%%%%%%%%%%%%%%%%%%%%%%%%%%%%%%%%%% M %
% % 
% \ProcedurePrintDefault{}
\MassnahmenNG{Spezielle Maßnahmen}{Abort}{\label{m:abort}}{
\begin{itemize}
  \item Lagerung: sterile Vorlage und Lagerung nach
  Fritsch
  % TODO Querverweis zu Lagerungsarten einfügen
  \item \TxBeiVit: \TxMassVitMKBes
  \begin{itemize}
    \item Lagerung: Situationsgerecht
  \end{itemize}
  \item Abgegangene Gewebeteile sammeln
  \item Transportentscheidung: Abt. f. Gynäkologie und
  Geburtshilfe
\end{itemize}
} 




% % M %%%%%%%%%%%%%%%%%%%%%%%%%%%%%%%%%%%%%%%%%%%%%%%%%%% M %
% % 
% \ProcedurePrintDefault{}
\MassnahmenNG{Spezielle Maßnahmen}{Eileiterschwangerschaft,
Verdacht}{\label{m:eileiterschwangerschaft-verdacht}}{
\begin{itemize}
    \item Transportentscheidung: Abt. f. Gynäkologie
\end{itemize}
}




% % M %%%%%%%%%%%%%%%%%%%%%%%%%%%%%%%%%%%%%%%%%%%%%%%%%%% M %
% % 
% \ProcedurePrintDefault{}
\MassnahmenNG{Spezielle Maßnahmen}{Eileiterschwangerschaft,
rupturiert}{\label{m:eileiterschwangerschaft-rupturiert}}{%
\begin{itemize}
  \item \TxMassVitMKBes
  \begin{itemize}
    \item Lagerung: Bauchdeckenentspannend.
    \item Schockbekämpfung
  \end{itemize}
\end{itemize}
}




% % M %%%%%%%%%%%%%%%%%%%%%%%%%%%%%%%%%%%%%%%%%%%%%%%%%%% M %
% % 
% \ProcedurePrintDefault{}
\MassnahmenNG{Spezielle Maßnahmen}{Vorzeitige
Plazentalösung}{\label{m:plazentaloesung}}{
\begin{itemize}
  \item \TxMassVitMKBes
  \begin{itemize}
    \item Lagerung: sterile Vorlage und Lagerung nach Fritsch.
    \item Schockbekämpfung
  \end{itemize}
  \item Keine vaginale Untersuchung oder  Scheidentamponade
  \item Zügiger Transport mit Voranmeldung
\end{itemize}
}




% % M %%%%%%%%%%%%%%%%%%%%%%%%%%%%%%%%%%%%%%%%%%%%%%%%%%% M %
% % 
% \ProcedurePrintDefault{}
\MassnahmenNG{Spezielle
Maßnahmen}{Vena-cava-Syndrom}{\label{m:vena-cava-syndrom}}{

\begin{itemize}
    \item Linksseitenlage, grundsätzlich immer bei
    hochschwangeren Patientinnen.
\end{itemize}
}




% % M %%%%%%%%%%%%%%%%%%%%%%%%%%%%%%%%%%%%%%%%%%%%%%%%%%% M %
% % 
% \ProcedurePrintDefault{}
\MassnahmenNG{Spezielle Maßnahmen}{Verdacht auf
Präeklampsie}{\label{m:praeeklampsie-verdacht}}
{
\begin{itemize}
    \item schonender Transport
    \item Überwachung, evtl. \ce{O2}
    \item Reizabschirmung
    %\item Notarzt: Krampfdurchbrechung, RR-Senkung
    %\item (im Vollbild: vorzeitige Beendigung der  Schwangerschaft)
\end{itemize}
}




% % M %%%%%%%%%%%%%%%%%%%%%%%%%%%%%%%%%%%%%%%%%%%%%%%%%%% M %
% % 
% \ProcedurePrintDefault{}
\MassnahmenNG{Spezielle Maßnahmen}{Eklampsie --
Eklamptischer Krampfanfall}{\label{m:eklampsie}}{
Grundsätzlich ist so wie bei jedem anderen Krampfanfall zu verfahren , vgl. Kap. \ref{AASS-topic:krampfanfall}. Bedenke:
Der eklamptische Krampfanfall ist für Mutter und Kind lebensgefährlich!
}




% % M %%%%%%%%%%%%%%%%%%%%%%%%%%%%%%%%%%%%%%%%%%%%%%%%%%% M %
% % 
% \ProcedurePrintDefault{}
\MassnahmenNG{Spezielle Maßnahmen}{Vorzeitiger
Fruchtwasserabgang}{\label{m:fruchtwasserabgang}}{
\begin{itemize}
    \item Lagerung: Linksseitenlagerung, Becken erhöht
    (Schwerkraft verzögert Geburt), Lagerung nach Fritsch
    \item Nicht mehr aufstehen lassen! Liegender Transport!
    \item Transportentscheidung: Abteilung für Gynäkologie und
    Geburtshilfe, Kreißsaal
\end{itemize}
}





% % M %%%%%%%%%%%%%%%%%%%%%%%%%%%%%%%%%%%%%%%%%%%%%%%%%%% M %
% % 
% \ProcedurePrintDefault{}
\MassnahmenNG{Spezielle Maßnahmen}{Geburt}{\label{m:geburt}}{
Vorbereitung:
\begin{itemize}
  \item (Not-)Arzt hinzuziehen
  \item Fahrzeug/Umgebung einheizen
  \item Vorbereiten des Geburtenset
  \item Lagerung: Z.\,B. erhöhter Oberkörper, angewinkelte
  Beine; Wunsch der Mutter berücksichtigen!
\end{itemize}
Geburt:
\begin{itemize}
  \item Stuhlschutz
  \item Versorgung des Neugeborenen
  (\prettyref{m:neugeborenes})
  \item Nachdem die Nabelschnur durchgeschnitten
  wurde, mütterliche Seite an
  Oberschenkel fixieren und vor Zug schützen
  \item Dokumentation von Geburtszeit und Geburtsort
\end{itemize}
}





% % M %%%%%%%%%%%%%%%%%%%%%%%%%%%%%%%%%%%%%%%%%%%%%%%%%%% M %
% % 
% \ProcedurePrintDefault{}
\MassnahmenWide{Spezielle Maßnahmen}
{Versorgung des Neugeborenen}
{\label{m:neugeborenes}}
{
\begin{itemize}
  \item Neugeborenes abrubbeln
  \item Erstbeurteilung:
  \begin{itemize}
    \item Muskeltonus
    \item Atemweg
    \begin{itemize}
      \item Bei Atemwegsverlegung: Mittels Oro-Sauger absaugen
      \hfill\Commentary[Absaugen Neugeborener]{ASBÖ
      Lehrmeinung 18/2011 \enquote{Absaugen Neugeborener im
      Rettungsdienst} \cite{AsboeLehrmeinung2011-19}: \Speech{\enquote{Mehrere
      Studien zeigten keine Verbesserung im Outcome
      gegenüber abgesaugter Neugeborener. Daher ist es
      obligat, keine Verzögerung bei den BLS-Maßnahmen
      hervorzurufen, welche durch lange Absaugversuche
      zustande kommen können. Lediglich die
      Atemwegsobstruktion durch Mekonium muss behandelt
      werden, indem die Atemwege freigemacht werden. Selbst
      missfarbiges Fruchtwasser stellt keine Indikation zur
      Absaugung
      dar.}}
      \ParBugfix
      ERC 2010 \cite{Erc2010Sek06De}: \Speech{\enquote{Absaugen ist
      nur notwendig, wenn die Atemwege verlegt sind. Eine
      solche Verle- gung kann aufgrund von Mekonium, auch
      wenn das Neugeborene keine Mekoniumablagerungen auf
      der Haut zeigt, aber auch Blutkoageln, zähem Schleim
      oder Vernix bestehen. Wird abgesaugt, ist zu bedenken,
      dass zu heftiges oropharyngeales Ab- saugen das
      Einsetzen einer suffizienten Spontanatmung verzögern
      und zu einem Laryngospasmus sowie zu einer vagalen
      Bradykardie führen kann. Das Vorhandensein von
      zähflüssigem Mekonium beim schlaffen, avitalen
      Neugeborenen ist die einzige Situation, in der ein
      sofortiges Absaugen des Oropharynx zu erwägen ist.}}
      }%Commentary
    \end{itemize}
    \item Atmung
    
    Bei Schnappatmung oder fehlender Atmung: Atemwege öffnen
    5 initiale Beatmungen
    \item Kreislauf: \Abk{HF} messen (durch Auskultation
    an der Herzspitze (nur geübtes Personal) oder an
    \E{Oberarmarterie}). Wenn ein \E{geeignetes
    Pulsoxymeter} vorhanden ist, kann auch dieses verwendet werden.
    \hfill\Commentary[Ort der HF-Messung beim Neugeborenen]
    {ERC 2010 \cite{Erc2010Sek06De}: \Speech{\enquote{Die beste Methode
    zur Beurteilung der Herzfrequenz ist die direkte Auskultation mit dem Stethoskop über der Herzspitze. Das Tasten des
    Pulses an der Basis der Nabelschnur ist oft möglich, kann
    aber irreführend sein. Eine Beurteilung der Herzfrequenz
    allein über die Pulsation der Nabelschnur ist nur
    zuverlässig, wenn die Herzfrequenz über 100\PerMin* liegt.}}
    \ParBugfix
    Bei Sanitätern, insbesondere der
    niedrigeren Ausbildungsstufen, kann nicht davon
    ausgegangen werden, dass eine Auskultation in dieser
    stressreichen Situation fehlerfrei und zuverlässig angewendet werden kann. Die
    Oberarmarterie erscheint in diesem Fall als
    pragmatischer Kompromiss.}
  \end{itemize}
\end{itemize}
{\FormatTable
\begin{tabularx}{\linewidth}{XX}
\toprule\addlinespace
\multicolumn{1}{c}{\includegraphics[width=1cm]{./icons/sm-basic.png}}
&
\multicolumn{1}{c}{\includegraphics[width=1cm]{./icons/sm-pale.png}~~\includegraphics[width=1cm]{./icons/sm-pale-frown.png}}
\\\addlinespace
\noindent\TabHeading{Kräftiges Schreien/suffiziente Atmung}
&
\noindent\TabHeading{Insuffiziente Spontanatmung,
Atemstillstand oder Bradykardie}
\\\addlinespace\cmidrule(lr){1-1}\cmidrule(lr){2-2}\addlinespace
\begin{itemize}
  \item[\eye] Guter Muskeltonus,
  \item[\eye] Herzfrequenz \textgreater\,100\PerMin*
\end{itemize}




&
\begin{itemize}
  \item[\eye] Normaler, reduzierter Muskeltonus, oder sogar
  schlaffer Muskeltonus (\enquote{floppy})
  \item[\eye] Herzfrequenz \textless\,100\PerMin*; Bradykarde
  oder nichtnachweisbare Herzfrequenz,
  \item[\eye] Oft mit ausgeprägter Blässe als Zeichen
  einer schlechten Perfusion.
\end{itemize}

\\\addlinespace\cmidrule(lr){1-1}\cmidrule(lr){2-2}\addlinespace
\begin{itemize}
  \item Abtrocknen und Einwickeln in warme Tücher:Warm
  einpacken (nur das Gesicht darf frei bleiben) und der Mutter
  überreichen.
  \item Abnabeln, frühestens nach 1\Min*:
  \begin{enumerate}
    \item Erst wird eine Handbreite vom Kind
    entfernt die erste Nabelklemme gesetzt.
    \item Dann wird die Nabelschnur Richtung Mutter ausgestreift und
    wieder eine Handbreit entfernt die zweite Klemme gesetzt.
    \item Zwischen den beiden
    Klemmen wird die Nabelschnur durchtrennt.
  \end{enumerate}
  \item Das Neugeborene
  kann der Mutter auf den Bauch gelegt werden. Durch den
  Kontakt zur Haut der zugedeckten Mutter wird das Baby
  gewärmt.
  
  \item Es kann zu diesem Zeitpunkt bereits an die Brust
  angelegt werden.
\end{itemize}
  &
  % Nachdem ein Neugeborenes dieser Gruppe getrocknet und in
  % warme Tücher gewickelt wurde, ist meist eine kurze
  % Maskenbeatmung ausreichend. Kommt es darunter nicht zu einem
  % adäquaten Anstieg der Herzfrequenz, müssen Herzdruckmassagen
  % durchgeführt werden.
  
  \begin{itemize}
    %     \item \ce{O2}-Gabe (\enquote{\ce{O2}-Dusche}) für 30\,Sek.
    \item Atemwege freimachen, absaugen
    \item 5 Beatmungen
    \item Wenn keine Steigerung der Herzfrequenz:
    \begin{itemize}
      %     \item \ce{O2}-Gabe (\enquote{\ce{O2}-Dusche}) für 30\,Sek.
      \item Lagerung und Atemwege überprüfen
      \item 5 Beatmungen wiederholen
    \end{itemize}
    \item Wenn weiterhin keine ausreichende Atmung oder
    \EE{Herzfrequenz \textless\,60\PerMin*}:
    \begin{itemize}
      \item \Sign{→} Reanimation des Neugeborenen
    \end{itemize}
    \FullPrettyRef{m:neugeborenes-basisreanimation}
  \end{itemize}
  
  % Nachdem ein Neugeborenes dieser Gruppe getrocknet und in
  % warme Tücher gewickelt wurde, müssen unverzüglich die
  % Atemwege geöffnet werden. Die Lungen müssen belüftet werden,
  % und das Kind muss beatmet werden. Wurden die Atemwege
  % erfolgreich geöffnet und die Lungen belüftet, können
  % außerdem Herzdruckmassagen und möglicherweise eine
  % Medikamentengabe notwendig sein. Eine sehr kleine Gruppe von
  % Neugeborenen bleibt trotz adäquater Spontanatmung und guter
  % Herzfrequenz hypox­ ämisch. Die Diagnosen umfassen hier ein
  % weites Spektrum, wie z. B. eine Zwerchfellhernie,
  % Surfactant-Mangel, eine kongenitale Pneumonie, einen
  % Pneumothorax oder ein angeborenes zyanotisches Herzvitium.
  
  
  \\\addlinespace\bottomrule
\end{tabularx}
}

\cite{Erc2010Sek06De,AsboeLehrmeinungRd2011-19,AsboeLehrmeinungRd2011-18}
}




% % M %%%%%%%%%%%%%%%%%%%%%%%%%%%%%%%%%%%%%%%%%%%%%%%%%%% M %
% % 
% \ProcedurePrintDefault{}
\MassnahmenNG{Spezielle Maßnahmen}
{Basisreanimation des Neugeborenen}
{\label{m:neugeborenes-basisreanimation}}
{
\begin{itemize}
  \item 2 Daumen nebeneinander über dem
  unteren Drittel des Brustbeins, direkt unter einer gedachten Linie zwischen den
  Brustwarzen
  
  \item Umgreifen des gesamten
  Brustkorbs und Unterstützung des Rückens des Kindes.
  \item Thoraxkompressionen und Beatmungen im Verhältnis von
  \EE{3\,:\,1}
  
%   \item Ziel: 90 Kompressionen und 30 Beatmungen\PerMin*
%   (120 Maßnahmen\PerMin*)
  \item \TxMassVitMK
\end{itemize}

\cite{Erc2010Sek06De}
}






% % M %%%%%%%%%%%%%%%%%%%%%%%%%%%%%%%%%%%%%%%%%%%%%%%%%%% M %
% % 
% \ProcedurePrintDefault{}
\MassnahmenNG{Spezielle Maßnahmen}{Nabelschnurvorfall}
{\label{m:nabelschnurvorfall}}
{
\begin{itemize}
  \item \Ca{Zeitkritisch!}
    \item \Ca{Notarztnachforderung nur wenn Stabilisierung
    notwendig}
	\item Beckenhochlagerung
	\item Mutter soll \underline{nicht} pressen (wenn möglich Geburt erst im
	Krankenhaus)
	\item Wehenhemmung (hecheln; \Z{medikamentös durch Notarzt})
	\item Nabelschnurumschlingung des kindlichen Halses während der Geburt sofort lösen
\end{itemize}
}






% % M %%%%%%%%%%%%%%%%%%%%%%%%%%%%%%%%%%%%%%%%%%%%%%%%%%% M %
% % 
% \ProcedurePrintDefault{}
\MassnahmenNG{Spezielle Maßnahmen}{Pathologische
Geburtslagen}
{\label{m:pathologische-geburtslagen}}
{
\begin{itemize}
  \item Bei Geburtsunmöglichen Lagen:\begin{itemize}
    \item \Ca{Zeitkritisch!}
    \item \Ca{Notarztnachforderung nur wenn Stabilisierung
    notwendig}
    \item Voranmeldung!
  \end{itemize}
	\item Beckenhochlagerung
	\item Mutter soll \underline{nicht} pressen (wenn möglich Geburt erst im
	Krankenhaus)
	\item Wehenhemmung (hecheln; \Z{medikamentös durch Notarzt})
\end{itemize}
}





% % M %%%%%%%%%%%%%%%%%%%%%%%%%%%%%%%%%%%%%%%%%%%%%%%%%%% M %
% % 
% \ProcedurePrintDefault{}
\MassnahmenNG{Spezielle Maßnahmen}{Placenta praevia und bevorstehende
Geburt}{\label{m:placenta-praevia}}
{
\begin{itemize}
    \item \TxMassVitMKBes
    \begin{itemize}
      \item \Ca{Zeitkritisch!}
      \item \Ca{Notarztnachforderung nur wenn Stabilisierung
      notwendig}
        \item Lagerung: Sterile Vorlage und Lagerung nach
        Fritsch.
        \item Ggfs. Schockbekämpfung
    \end{itemize}
\end{itemize} \autocite{MA70-Protokoll-2012-10-11}
}




% % M %%%%%%%%%%%%%%%%%%%%%%%%%%%%%%%%%%%%%%%%%%%%%%%%%%% M %
% % 
% \ProcedurePrintDefault{}
\MassnahmenNG{Spezielle
Maßnahmen}{Uterusatonie}{\label{m:uterusatonie}}{
\begin{itemize}
  \item \TxMassVitMKBes
  \begin{itemize}
    \item Schockbekämpfung
    \item Lagerung: Fritsch, Beine hoch
  \end{itemize}
\end{itemize}
}




% % M %%%%%%%%%%%%%%%%%%%%%%%%%%%%%%%%%%%%%%%%%%%%%%%%%%% M %
% % 
% \ProcedurePrintDefault{}
\MassnahmenNG{Spezielle
Maßnahmen}{Uterusruptur}{\label{m:uterusruptur}}{
\begin{itemize}
	\item \TxMassVitMK
	\item Schockbekämpfung
	\item Lagerung: Fritsch, Beine hoch
\end{itemize}
}




% % M %%%%%%%%%%%%%%%%%%%%%%%%%%%%%%%%%%%%%%%%%%%%%%%%%%% M %
% % 
% \ProcedurePrintDefault{}
\MassnahmenNG{Spezielle Maßnahmen}{Asphyxie des
Neugeborenen}
{\label{m:asphyxie}}
{%
Die Behandlung der Asphyxie erfolgt im Rahmen der Versorgung
des Neugeborenen (\FullPrettyRef{m:neugeborenes}).
% \begin{itemize}
%   \item \TxMassVitMK
%   \item Warmhalten
%   \item Abtrocknen/abrubbeln; stimulieren: Fußsohlen und
%   Rücken klopfen, reiben
%   \item Atmung und Herzfrequenz beurteilen
%   \item Bei Schnappatmung, fehlender Atmung oder HF <
%   100\PerMin*:  \begin{itemize}
% %     \item \ce{O2}-Gabe (\enquote{\ce{O2}-Dusche}) für 30\,Sek.
%     \item Atemwege freimachen, absaugen
%     \item 5 Beatmungen
%   \end{itemize}
%   \item Wenn keine Steigerung der Herzfrequenz: 
%   \begin{itemize}
% %     \item \ce{O2}-Gabe (\enquote{\ce{O2}-Dusche}) für 30\,Sek.
%     \item Lagerung und Atemwege überprüfen 
%     \item 5 Beatmungen wiederholen
%     \item \ce{O2}-Gabe
%     \end{itemize}
%   \item Wenn weiterhin keine ausreichende Atmung oder
%   \EE{Herzfrequenz < 60\PerMin*} 
%   
%   \Sign{→} Reanimation des Neugeborenen
% \end{itemize}
}





% % M %%%%%%%%%%%%%%%%%%%%%%%%%%%%%%%%%%%%%%%%%%%%%%%%%%% M %
% % 
% \ProcedurePrintDefault{}
\MassnahmenNG{Spezielle Maßnahmen}{Vaginale
Blutung}{\label{m:vaginale-blutung}}{
\begin{itemize}
    \item Lagerung nach Fritsch
    \item Transportentscheidung: Abteilung für
    Gynäkologie, bei Wehen Kreißsaal
\end{itemize}
}




% % M %%%%%%%%%%%%%%%%%%%%%%%%%%%%%%%%%%%%%%%%%%%%%%%%%%% M %
% % 
% \ProcedurePrintDefault{}
\MassnahmenNG{Spezielle
Maßnahmen}{Laryngitis}{\label{m:laryngitis}}{

\begin{itemize}
    \item Eltern und Kind beruhigen.
    \item Feuchte Luft einatmen lassen (evtl. Fenster
    öffnen; feuchte Tücher im Raum aufhängen; Dusche im Badezimmer aufdrehen)
\end{itemize}
}




% % M %%%%%%%%%%%%%%%%%%%%%%%%%%%%%%%%%%%%%%%%%%%%%%%%%%% M %
% % 
% \ProcedurePrintDefault{}
\MassnahmenNG{Spezielle
Maßnahmen}{Epiglottitis}{\label{m:epiglottitis}}{%

\begin{itemize}
    \item \TxMassVitMK
    \item Luftbefeuchtung
    \item Kinder sitzen lassen!
    \item Keine Manipulation im Mundraum (Spatel)!
\end{itemize}
}




% % M %%%%%%%%%%%%%%%%%%%%%%%%%%%%%%%%%%%%%%%%%%%%%%%%%%% M %
% % 
% \ProcedurePrintDefault{}
\MassnahmenNG{Spezielle Maßnahmen}{SIDS}{\label{m:sids}}{%

\begin{itemize}
	\item Reanimation.
	\item Betreuung der Angehörigen. Den Eltern muss das Gefühl
	gegeben werden, dass sie alles für ihr Kind getan haben. Ein solcher
	Unglücksfall tritt schicksalshaft auf und wäre nicht abwendbar gewesen. 
	\item In jedem Fall nimmt die Kriminalpolizei Ermittlungen auf. Darauf müssen
	die Eltern vorbereitet werden. 
	\item Betreuung für die Angehörigen organisieren
	(Kriseninterventionsdienst\opt{REGVIE}{\footnote{In Wien ist die
	\E{Akutbetreuung Wien} der Stadt Wien zuständig. Sie kann über die
	Leitstelle angefordert werden.}})
	\item Zwillingsgeschwister unbedingt hospitalisieren.
	\item Einsatznachbesprechung und eigene Psychohygiene. Derartige Einsätze
	können -- auch für erfahrenes Fachpersonal -- sehr belastend sein. 
	
    Meldung an Gesamteinsatzleiter, dienstführenden
    Offizier o.\,ä., Erkundigung nach Verfügbarkeit von
    Peers (auch wenn Peer-Angebot nicht angenommen wird!). 
\end{itemize}

}




% % M %%%%%%%%%%%%%%%%%%%%%%%%%%%%%%%%%%%%%%%%%%%%%%%%%%% M %
% % 
% \ProcedurePrintDefault{}
\MassnahmenNG{Spezielle
Maßnahmen}{Ertrinkungsunfall}{\label{m:ertrinkungsunfall}}{
\begin{itemize}
  \item Wenn erforderlich Reanimation mit fünf
  Initialbeatmungen. Die Reanimation kann auch nach länger
  andauerndem Kreislaufstillstand erfolgreich sein, da bei
  einer Unterkühlung der Stoffwechsel verlangsamt wird.
  \item \TxBeiVit \TxMassVitMK
  \item Bei Aspiration: Wenn der Verdacht auf eine
  Aspiration besteht, ist der Patient zur Beobachtung zu
  hospitalisieren.
\end{itemize}
}




% % M %%%%%%%%%%%%%%%%%%%%%%%%%%%%%%%%%%%%%%%%%%%%%%%%%%% M %
% % 
% \ProcedurePrintDefault{}
\MassnahmenNG{Spezielle Maßnahmen}{Fieberkrampf im
Kindesalter}{\label{m:fieberkrampf-kindesalter}}{
    
    Die Behandlung erfolgt grundsätzlich wie beim Erwachsenen (\FullPrettyRef{m:zerebraler-krampfanfall-bestendend}).
    Die Körpertemperatur ist immer zu messen! 
}





% % M %%%%%%%%%%%%%%%%%%%%%%%%%%%%%%%%%%%%%%%%%%%%%%%%%%% M %
% % 
% \ProcedurePrintDefault{}
\MassnahmenNG{Spezielle Maßnahmen}
{Vergewaltigung und geschlechtliche Nötigung}
{\label{m:vergewaltigung}}
{
\begin{itemize}
  \item Ruhig und sachlich bleiben
  \item Polizei verständigen
  \item Ermöglichung von polizeilichen Ermittlungen unter Schonung
  des Opfers (\enquote{vermitteln})
  \item Besonders wichtig: Ausführliche sachliche Erklärung
  aller durchzuführenden Maßnahmen
  \item Symptomatische Betreuung von Begleitverletzungen
  \item Keine übertriebenen Untersuchungen
  \item Keine gynäkologischen Untersuchungen
  \item Transportentscheidung: Abt. f. Gynäkologie
\end{itemize}

}




% % M %%%%%%%%%%%%%%%%%%%%%%%%%%%%%%%%%%%%%%%%%%%%%%%%%%% M %
% % 
% \ProcedurePrintDefault{}
\MassnahmenNG{Spezielle
Maßnahmen}{Kindesmisshandlung}{\label{m:kindesmisshandlung}}{
Die Aufgaben des Rettungsdienstes bei Verdacht auf Kindesmisshandlung besteht in:

\begin{enumerate}
  \item Daran denken!
  \item Umstände und Umfeld abklären
  \item Unstimmigkeiten und Auffälligkeiten auf eine neutrale Weise
  \EE{dokumentieren}
  \item Obiges verlässlich bei der Übergabe an die nachbetreuende Einrichtung
  weiterleiten. Evtl. Name des Gesprächpartners bei der  Übergabe dokumentieren.
  \item \EE{Patienten nicht im Stich lassen}: Gelindestes Mittel wählen,
  evtl. Vorwand für die Hospitalisierung suchen, aber wenn ein begründeter
  Verdacht auf eine Misshandlung besteht und die
  Erziehungsberechtigten unkooperativ sind notfalls auch die Exekutive beiziehen.
\end{enumerate}

\paragraph{Vorsicht!}

Jedenfalls zu \underline{unterlassen} sind:

\begin{itemize}
  \item Voreilige Verdachtsäußerungen. \Z{Nicht alles was stinkt ist ein
  Fisch!}
  \item Andeutungen gegenüber den Erziehungsberechtigten oder Dritten.
\end{itemize}

}




% % M %%%%%%%%%%%%%%%%%%%%%%%%%%%%%%%%%%%%%%%%%%%%%%%%%%% M %
% % 
% \ProcedurePrintDefault{}
\MassnahmenNG{Spezielle Maßnahmen}{Umgang mit Patienten mit
Wahnvorstellungen}{\label{m:wahnvorstellungen}}{ Die -- für
den Außenstehenden -- wirren Vorstellungen stellen das Fachpersonal vor die Frage, wie mit diesen
Patienten umzugehen ist.
Ziel ist:
\begin{enumerate}
  \item Abwendung von unmittelbaren Gefahren für den
  Patienten und die Umgebung
  \item Korrekte und professionelle Durchführung eines
  Transportes in eine geeignete Einrichtung.
\end{enumerate}
Es gelten somit folgende Grundsätze:
\begin{enumerate}
  \item Sofern eine Gefährdung für den Patienten oder für
  Dritte besteht, ist diese (wenn möglich) zu beseitigen,
  dazu zählt u. U. auch der Rückzug und die Herbeiziehung
  der Exekutive.
  \item Wahnvorstellungen sollen \EE{nicht ausgeredet}
  werden.
  \item  Wahnvorstellungen sollen aber auch \EE{nicht
  bestätigt} werden.
  \item Am besten man nimmt die Schilderungen des
  Patienten \EE{interessiert, aber neutral und möglichst
  kommentarlos} zur Kenntnis (\emph{\enquote{Mhm!}, \enquote{Aha.}}).
\end{enumerate}
}




% % M %%%%%%%%%%%%%%%%%%%%%%%%%%%%%%%%%%%%%%%%%%%%%%%%%%% M %
% % 
% \ProcedurePrintDefault{}
\MassnahmenNG{Spezielle Maßnahmen}{Unruhiger oder
aggressiver Patient}{\label{m:aggressiver-patient}}{

\begin{itemize}
  \item Ruhe vermitteln und deeskalieren:
  \begin{itemize}
  \item  Ruhig bleiben! Achtung auf eigenen Tonfall!
  \item Nicht provozieren!
  \item Aufregung nicht abfärben lassen
  \item \emph{\enquote{Talk down}}
  \end{itemize}
  \item Selbstschutz, Fluchtweg offen halten
\end{itemize}

}




% % M %%%%%%%%%%%%%%%%%%%%%%%%%%%%%%%%%%%%%%%%%%%%%%%%%%% M %
% % 
% \ProcedurePrintDefault{}
\MassnahmenWide{Spezielle Maßnahmen}{Hitzekollaps,
Hitzeerschöpfung,
Hitzschlag}{\label{m:systemische-hitzeschaedigungen}}{
\FormatTable

\begin{tabularx}{\linewidth}{XXX}
\TabHeading{Hitzekollaps}
&
\TabHeading{Hitzerschöpfung}
&
\TabHeading{Hitzschlag}
\\\addlinespace\midrule\addlinespace
% \begin{minipage}{\linewidth}
\begin{itemize}
  \item Flachlagerung: ${\rightarrow}$ Gehirndurchblutung
  wieder ${\uparrow}$, bewusstseinsklar
  \item nach Sekunden gebessert, \enquote{regelt sich von
  selbst}
  \item Patient sitzen lassen, Flüssigkeit
\end{itemize}
% \end{minipage}
&
% \begin{minipage}{\linewidth}
\begin{itemize}
  \item Kühlung, Schatten
  \item Flüssigkeitszufuhr, wenn bewusstseinsklar
  \item Neurocheck inkl. BZ-Messung
\end{itemize}
% \end{minipage}
&
% \begin{minipage}{\linewidth}
\begin{itemize}
  \item \TxMassVitMK
  \item Beengende Kleidungsstücke lockern
  \item Flachlagerung an kühlem Ort (Schatten), Beine und Kopf hochlagern. 
  \item Kühlung von \E{außen} (Fächer, mit Eiswürfeln abreiben etc.)
  \item Kühlung von \E{innen}: für \EE{ansprechbare} Patienten viel kühle
  (alkoholfreie) Flüssigkeit, evtl. Elektrolyt-Getränke
  \item Neurocheck inkl. BZ-Messung
\end{itemize}
% \end{minipage}
\\\addlinespace\bottomrule
\end{tabularx}
}




% % M %%%%%%%%%%%%%%%%%%%%%%%%%%%%%%%%%%%%%%%%%%%%%%%%%%% M %
% % 
% \ProcedurePrintDefault{}
\MassnahmenNG{Spezielle
Maßnahmen}{Sonnenstich}{\label{m:sonnenstich}}{

\begin{itemize}
  \item Vitale Bedrohung einschätzen
  \TxBeiVit \TxMassVitMK
  \item Flachlagerung an kühlem Ort (Schatten), Kopf hochlagern
  \item Beengende Kleidungsstücke lockern
  \item Kühlung von außen (Fächer, mit Eiswürfeln abreiben, kalte
  Umschläge auf die Stirn etc.)
  \item Kühlung von innen: für \EE{ansprechbare} Patienten kühle
  (alkoholfreie) Flüssigkeit, evtl. Elektrolyt-Getränke
%   \item Schocklagerung u. -bekämpfung. % REVIEW 2010-05-19
  \item Neurocheck inkl. BZ-Messung
\end{itemize}
}




% % M %%%%%%%%%%%%%%%%%%%%%%%%%%%%%%%%%%%%%%%%%%%%%%%%%%% M %
% % 
% \ProcedurePrintDefault{}
\MassnahmenNG{Spezielle Maßnahmen}{Schwere Unterkühlung
(\textless\,34\DegC*)}{\label{m:unterkuehlung-schwer}}{

\begin{itemize}
  \item \TxMassVitMK
  \item Den Patient so wenig als möglich (nur soviel wie unbedingt nötig)
  bewegen (drohender Bergungstod).
  \item Schutz vor weiterer Abkühlung, \Ca{keine \underbar{aktive}
  Erwärmung!}
  \begin{itemize}
    \item Woll- und Aludecken! Mütze!
  \end{itemize}
\end{itemize}
}




% % M %%%%%%%%%%%%%%%%%%%%%%%%%%%%%%%%%%%%%%%%%%%%%%%%%%% M %
% % 
% \ProcedurePrintDefault{}
\MassnahmenWide{Allgemeine Maßnahmen und
Herangehensweise}{Vergiftungen}{\label{m:am-vergiftung}}{
\begin{enumerate}
  \item Selbstschutz
  \item Beendigung des Giftkontaktes
  \item Evtl. Dekontamination bei
  Gefahrstoffen\footnote{z.\,B. Insektizide, Lösungsmittel
  in der chemischen Industrie, \ldots}
  \item Sicherung der Vitalfunktionen
  \item Entscheidung, ob es sich um einen akut vital
  bedrohten Patienten handelt
  \item Quellenforschung (Mistkübel
  durchsuchen)
  \item Nüchtern lassen
  \item NICHT absichtlich Erbrechen
  herbeiführen
  \item Asservierung (Packungen, Erbrochenes in
  Nierentasse mitnehmen!)
  \item Anamnese, Besonderheiten:
  \begin{multicols}{2}
\begin{enumerate}\raggedbottom
  \item \EE{Wer}? (Alter, KG, Vorerkrankungen, \ldots
  $\rightarrow$ Risikoabschätzung!)
  \item \EE{Was}? (Art des Giftes $\rightarrow$ evt.
  zu erwartende Symptome)
  \item \EE{Wann}? (Zeit seit Einnahme)
  \item \EE{Wie}?
  (inhalativ, oral, perkutan, intravenös etc.)
  \item \EE{Wieviel}? (Dosisabschätzung)
  \item \EE{Warum}? (Ursachenforschung:
  Selbst\-mord\-ver\-such, Unfall, Fremdverschulden,
  Vorsatz?)
\end{enumerate}
\end{multicols}
\item Ggfs. \EE{Expertenrat einholen} und Entscheidung
bezüglich des weiteren Vorgehens

\begin{multicols}{2}
\begin{enumerate}\raggedbottom
  \item Gegenmittel verfügbar?
  \item \enquote{Normales} Spitalsbett?
  \item Intensivüberwachung?
  \item Intensivüberwachung mit Entgiftungseinheit?
  \item Dialyse?
  \item Druckkammer?
  \item Sonstige Spezialbehandlungseinheit?
\end{enumerate}
\end{multicols}
\end{enumerate}
}






% % M %%%%%%%%%%%%%%%%%%%%%%%%%%%%%%%%%%%%%%%%%%%%%%%%%%% M %
% % 
% \ProcedurePrintDefault{}
\MassnahmenNG{Spezielle Maßnahmen}{Vergiftungen mit
Alkohol}{\label{m:intoxikation-alkohol}}{
\begin{itemize}
    \item Allgemeine Maßnahmen bei Vergiftungen
    (\prettyref{m:am-vergiftung})
    \item \TxBeiVit \TxMassVitMK
\end{itemize}
}




% % M %%%%%%%%%%%%%%%%%%%%%%%%%%%%%%%%%%%%%%%%%%%%%%%%%%% M %
% % 
% \ProcedurePrintDefault{}
\MassnahmenNG{Spezielle
Maßnahmen}{Opiatvergiftung}{\label{m:intoxikation-opiate}}{
\begin{itemize}
    \item Allgemeine Maßnahmen bei Vergiftungen
    (\prettyref{m:am-vergiftung})und
    symptomatische Therapie
    \item \TxBeiVit \TxMassVitMK
\end{itemize}

\Customized{
\begin{description}
  \item[\CustAsboe] Bei der Opiatintoxikation mit Atemstörung ist die Gabe
  von Naloxon nasal oder i.\,v. gem. Algorithmus durch NKA/NKV vorgesehen.
  \cite{Asb921Memo-AL-2011-000001}.
\end{description}
}

%Allgemeine Maßnahmen bei Vergiftungen
%(\ref{vergiftung-allgemein-massnahmen})und symptomatische
%Therapie. 
% Gegenmittel verfügbar.
}




% % M %%%%%%%%%%%%%%%%%%%%%%%%%%%%%%%%%%%%%%%%%%%%%%%%%%% M %
% % 
% \ProcedurePrintDefault{}
\MassnahmenNG{Spezielle Maßnahmen}{Vergiftung mit
Uppers}{\label{m:intoxikation-uppers}}{
\begin{itemize}
    \item Allgemeine Maßnahmen bei Vergiftungen
    (\prettyref{m:am-vergiftung}) und
    symptomatische Therapie
    \item \TxBeiVit \TxMassVitMK
\end{itemize}
}




% % M %%%%%%%%%%%%%%%%%%%%%%%%%%%%%%%%%%%%%%%%%%%%%%%%%%% M %
% % 
% \ProcedurePrintDefault{}
\MassnahmenNG{Allgemeine Maßnahmen}{Stickgasvergiftungen}{\label{m:am-stickgase}}{


\begin{enumerate}
  \item Auf Selbstschutz achten!
  \item Patienten retten, aber nur wenn es der
  Eigenschutz zulässt.
  \item \ce{O2}-Berieselung mit Maske ist in keinem Fall
  ein ausreichender Atemschutz!
  
  \item Spezialkräfte anfordern (Schwerer Atemschutz).
  I.\,d.\,R. Feuerwehr
  \item Weitere Opfer ausschließen (lassen):
  Nachbarwohnungen, etc.
  \item Sicherung der Vitalfunktionen, Beurteilung ob der
  Patient vital bedroht ist.
  \item \TxBeiVit \TxMassVitMKBes
  \begin{itemize}
    \item \ce{O2} hochdosiert, 15\LPerMin*
    \item Lagerung: situationsgerecht.
  \end{itemize}
  \item Traumacheck (Folgeverletzungen)
  \item Neurocheck inkl. Blutzuckerbestimmung
  (Differentialdiagnosen, Beurteilung der Beeinträchtigung im Verlauf)
\end{enumerate}
}




% % M %%%%%%%%%%%%%%%%%%%%%%%%%%%%%%%%%%%%%%%%%%%%%%%%%%% M %
% % 
% \ProcedurePrintDefault{}
\MassnahmenNG{Spezielle
Maßnahmen}{Kohlenmonoxid-Vergiftung}{\label{m:intoxikation-kohlenmonoxid}}{

\begin{itemize}
    \item Allgemeine Maßnahmen bei
Stickgas-Vergiftungen
(\prettyref{m:am-stickgase})
\item \Z{Hyperbare \ce{O2}-Therapie in der Druckkammer erwägen}
\end{itemize}
}




% % M %%%%%%%%%%%%%%%%%%%%%%%%%%%%%%%%%%%%%%%%%%%%%%%%%%% M %
% % 
% \ProcedurePrintDefault{}
\MassnahmenNG{Spezielle
Maßnahmen}{Kohlendioxid-Vergiftung}{\label{m:intoxikation-kohlendioxid}}{
Allgemeine Maßnahmen bei Stickgas-Vergiftungen
(\prettyref{m:am-stickgase})
}




% % M %%%%%%%%%%%%%%%%%%%%%%%%%%%%%%%%%%%%%%%%%%%%%%%%%%% M %
% % 
% \ProcedurePrintDefault{}
\MassnahmenNG{Allgemeine
Maßnahmen}{Reizgasvergiftungen}{\label{m:am-reizgase}}{


\begin{enumerate}
  \item Auf Selbstschutz achten!
  \begin{itemize}
    \item Patienten retten, aber nur wenn es der
    Eigenschutz zulässt.
    \item \ce{O2}-Berieselung mit Maske ist in keinem Fall
    ein ausreichender Atemschutz!
  \end{itemize}
  \item Gefahr erkennen: Welcher Stoff? Expertenrat einholen!
  \item Wenn erforderlich: Absperrmaßnahmen durchführen
  \item Wenn erforderlich:  Spezialkräfte anfordern
  \item Vorgehen analog zu Abschnitt \enquote{Gefahrengutunfall}
  (\prettyref{AASS-topic:gefahrengutunfall})
  \item Weitere Opfer ausschließen (lassen):
  Nachbarwohnungen, etc
  \item Sicherung der Vitalfunktionen, Beurteilung ob der
  Patient vital bedroht ist.
  \item \TxBeiVit 
  \item Traumacheck (Folgeverletzungen)
  \item Neurocheck inkl. Blutzuckerbestimmung
  (Differentialdiagnosen, Beurteilung der Beeinträchtigung im Verlauf)
\end{enumerate}
}




% % M %%%%%%%%%%%%%%%%%%%%%%%%%%%%%%%%%%%%%%%%%%%%%%%%%%% M %
% % 
% \ProcedurePrintDefault{}
\MassnahmenNG{Spezielle Maßnahmen}{Einnahme von Säuren oder
Laugen}{\label{m:saeuren-laugen}}{
\begin{itemize}
  \item Allgemeine Maßnahmen bei Vergiftungen
  (\prettyref{m:am-vergiftung})
  \item Schluckweise Wasser zur Verdünnung trinken lassen
  \item \EE{Niemals absichtlich Erbrechen  herbeiführen!}
  \item \Z{Keine Neutralisationsversuche! (Milch
  obsolet!)}
\end{itemize}
}




% % M %%%%%%%%%%%%%%%%%%%%%%%%%%%%%%%%%%%%%%%%%%%%%%%%%%% M %
% % 
% \ProcedurePrintDefault{}
\MassnahmenNG{Spezielle Maßnahmen}{Einnahme von
schaumbildenden Substanzen}{\label{m:schaumbildner}}{
\begin{itemize}
    \item Allgemeine Maßnahmen bei Vergiftungen
    (\prettyref{m:am-vergiftung})
    \item Patienten nüchtern lassen
    \item \EE{Niemals absichtlich Erbrechen 
    herbeiführen!}
\end{itemize}
}






% % M %%%%%%%%%%%%%%%%%%%%%%%%%%%%%%%%%%%%%%%%%%%%%%%%%%% M %
% % 
% \ProcedurePrintDefault{}
\MassnahmenNG{Allgemeine Maßnahmen}{Wundversorgung}
{\label{m:wundversorgung}}
{
\begin{itemize}
  \item Wunden dürfen nicht direkt berührt werden, außer es
  ist zur raschen Blutstillung unbedingt notwendig! 
Es
  muss immer mit Handschuhen und mit \EE{sterilem Material}
  gearbeitet werden!\footnote{Bei der Wundversorgung sind
  immer sterile Materialien zu verwenden. Papierhandtücher
  o.\,ä. sind \E{nicht zulässig}! Eine Ausnahme kann
  lediglich bei der Versorgung des Stichs bei der
  Blutzuckermessung mittels eines sauberen Tupfers gemacht
  werden, da die Wunde minimal klein ist.}
  
  % \item Wunden sollen grundsätzliche \E{nicht} gereinigt
  % werden! Sie dürfen nur lose Glassplitter oder ähnliches
  % vorsichtig aus der Wunde entfernen.
  \item \EE{Reinigung}: Oberflächliches \EE{ab}spülen mit
  steriler \I[Kochsalzlösung!pyhsiologische]{physiologischer
  Kochsalzlösung} (\ce{NaCl}
  0,9\PC*) von innen
  nach außen. Bei oberflächlichen
  Wunden ist auch ein Abspülen unter fließendem \IX{Leitungswasser} 
  zulässig \cite{Fernandez2012}. 
  Wunden, die im Spital
  weiter versorgt werden, sollen nicht übermäßig gereinigt werden, 
  auch soll keine
  wesentliche Zeitverzögerung entstehen.  
  \Commentary[Wundreinigung]{Alternativ sind
  statt der physiologischen Kochsalzlösung auch andere
  kristalloide Infusionslösungen verwendbar. Die
  Desinfektion mittels Wunddesinfektionsmittels ist nicht
  mehr empfohlen. \cite{Fernandez2012}}
  \item \EE{Fremdkörper} (Messer, Schere oder sonstige
  Pfählungsgegenstände) sind \EE{in der Wunde zu belassen}!
  Diese Gegenstände sind mit geeignetem Fixationsmaterial
  (z.\,B. Mullbinden) zu stabilisieren. Lose, kleine
  Fremdkörper wie z.\,B. Glassplitter dürfen mittels \E{steriler Instrumente} (Pinzette, {\ldots}) entfernt
  werden.
  \item Ausgetretene Strukturen (z.\,B. Gehirn, Eingeweide)
  dürfen \E{nicht zurückgestopft} werden! Solche Organteile
  sollten mit steriler physiologischer Kochsalzlösung
feucht gehalten und anschließend steril abgedeckt werden! 
  \item Ein steriler
  Wundverband mit der jeweiligen Situation angemessenem
  Material (z.\,B. Pflaster, Momentverband, Kompresse mit
  Mullbinde, Wundfolie, \ldots) muss angelegt oder die
  Wunde anderweitig steril, aber \E{nicht} luftdicht
  abgedeckt werden (Ausnahme: Bagatellverletzungen).
\end{itemize}
}
{}










% % M %%%%%%%%%%%%%%%%%%%%%%%%%%%%%%%%%%%%%%%%%%%%%%%%%%% M %
% % 
% \ProcedurePrintDefault{}
\MassnahmenNG{Allgemeine
Maßnahmen}{Unfälle}{\label{m:am-unfaelle}}{%

\begin{itemize}
  \item Unfallstelle absichern!
  \item Szeneüberblick: Wie viele Verletzte? Unfallszenario?
  \item Ggf. weitere Einsatzkräfte nachfordern!
  \item Beurteilung des Unfallmechanismus:
  \begin{itemize}
    \item Aus welcher Richtung hat die Kraft auf den Patienten eingewirkt?
    \item Wie schnell war das Fahrzeug unterwegs? bzw. Aus
    welcher Höhe ist der Patient gefallen?
    \item War der Patient angegurtet? Haben Airbags
    ausgelöst?
    \item Welche besondere Verformungen sind am Fahrzeug
    sichtbar?
    \item Wurde der Patient aus dem Fahrzeug geschleudert?
    bzw. Wurde ein Fußgänger erfasst?
  \end{itemize}
  \item Bei Hospitalisation: Patient \EE{nüchtern lassen!}
\end{itemize}
}




% % M %%%%%%%%%%%%%%%%%%%%%%%%%%%%%%%%%%%%%%%%%%%%%%%%%%% M %
% % 
% \ProcedurePrintDefault{}
\MassnahmenNG{Spezielle
Maßnahmen}{Fraktur}{\label{m:fraktur}}{
\EE{Taktik}: \Speech{\Ca{Zeitkritisch. Evtl. vitale
Bedrohung v.\,a. durch Blutverlust.}
Erkennen
von großen Blutverlusten, Blutstillung,  Steriles
Abdecken, Einschätzen der vitalen Bedrohung, ggfs
Schmerztherapie. 
Bei \E{vitaler Bedrohung}: Zeitkritisch,
zügiger Transport. 
Bei \E{stabilen Patienten}: Schienung}




\begin{itemize}
  \item  \EE{Blutstillung}: Z.\,B. mit Beckengurt
  (Beckenfraktur)
  \item Vitale Bedrohung einschätzen, ggfs. \TxMassVitMK.
  Ggfs. mit hypovolämischen Schock rechnen, Schockbekämpfung
  \item  \EE{Wundversorgung}: Bei offenen Frakturen ist
  Keimfreiheit be\-sonders wichtig, da  die Gefahr einer lebenslangen
  Knochen\-marks\-ent\-zündung \Z{(Osteomyelitis)} besteht!
  Wenn möglich soll eine \E{Wundfolie} (z.\,B.
  \index{OpSite}\T{OpSite}\texttrademark ) verwendet werden.
  \item Kleidung und Schmuck von den betroffenen Gliedmaßen
  entfernen (Schwellung!)
  \item Ggfs. Schmerztherapie (veranlassen)
  \item Ruhigstellung/Schienung bei stabilen Patienten
  \item Beurteilung von Durchblutung, Motorik und
  Sensibilität (DMS) an der betroffenen Extremität und an
  der Gegenseite
  \item Transportentscheidung: Abt. f. Unfallchirurgie, bei
  vitaler Bedrohung: Schockraum
\end{itemize}
}




% % M %%%%%%%%%%%%%%%%%%%%%%%%%%%%%%%%%%%%%%%%%%%%%%%%%%% M %
% % 
% \ProcedurePrintDefault{}
\MassnahmenNG{Spezielle
Maßnahmen}{Verstauchung}{\label{m:verstauchung}} {%
\begin{itemize}
  \item Kleidung und Schmuck von der betroffenen Gliedmaße entfernen
  (Vorsicht Schwellung!)
  \item Ruhigstellen
  \item Hochlagern
  \item Kühlen
  \item Beurteilung von Durchblutung, Motorik und
  Sensibilität (DMS) an der betroffenen Extremität und an
  der Gegenseite
  \item Transportentscheidung: Abt. f. Unfallchirurgie
%   (muss nicht sofort sein)
\end{itemize}
}




% % M %%%%%%%%%%%%%%%%%%%%%%%%%%%%%%%%%%%%%%%%%%%%%%%%%%% M %
% % 
% \ProcedurePrintDefault{}
\MassnahmenNG{Spezielle
Maßnahmen}{Verrenkung}{\label{m:verrenkung}}
{
\begin{itemize}
  \item Kleidung und Schmuck von der betroffenen Gliedmaße entfernen 
  (Vorsicht Schwellung!)
  \item Stellung der Gliedmaße belassen
  \item Federnd fixierte Lagerung
  \item \EE{Keine Einrenkversuche!}
  \item Beurteilung von Durchblutung, Motorik und
  Sensibilität (DMS) an der betroffenen Extremität und an
  der Gegenseite
  \item Transportentscheidung: Abt. f. Unfallchirurgie
\end{itemize}
}




% % M %%%%%%%%%%%%%%%%%%%%%%%%%%%%%%%%%%%%%%%%%%%%%%%%%%% M %
% % 
% \ProcedurePrintDefault{}
\MassnahmenNG{Spezielle Maßnahmen}{SHT}{\label{m:am-sht}\label{m:sht}}{

\begin{itemize}
  \item Vitale Bedrohung beurteilen. \par \TxBeiVit
  \TxMassVitMKBes
  \begin{itemize}
    \item Lagerung: s.\,u.
    \item Monitoring: Zusätzlich Beurteilung nach GCS (Verlauf beobachten und
    dokumentieren, \prettyref{AASS-topic:gcs})
  \end{itemize}
  \item Lagerung: leicht erhöhter Oberkörper, ca.
  30\,\textdegree{}, wenn keine Kontraindikation
  besteht.
  \item Je nach Unfallmechanismus \Abk{HWS}-Immobilisation (z.\,B. mit Stifneck\texttrademark)
  \item Ggfs. Wundversorgung
  \begin{itemize}
    \item Steriler Verband (steriles Abdecken)
    \item Ausgetretene Hirnmasse feucht steril abdecken
        und locker befestigen
    \item Knochen oder Fremdkörper \E{nicht} entfernen oder zurück stopfen
  \end{itemize}
  %        \begin{itemize}
%            \item Gedecktes \Abk{SHT}:
%            \begin{itemize}
%                \item Steriler Verband
%            \end{itemize}
%            \begin{itemize}
%                \item Steriler Verband locker auf offene
%                Verletzungen
%                \item Knochen / Hirnteile NICHT
%                zurückstopfen /  entfernen!
%                \item Ausgetretene Hirnmasse feucht steril abdecken
%            \end{itemize}
%        \end{itemize}
\end{itemize}
}




% % M %%%%%%%%%%%%%%%%%%%%%%%%%%%%%%%%%%%%%%%%%%%%%%%%%%% M %
% % 
% \ProcedurePrintDefault{}
\MassnahmenNG{Spezielle Maßnahmen}
{Rückenmarksverletzung, Verdacht}
{\label{m:ws-trauma-verdacht}}
{%
\EE{Taktik}: \Speech{
Immobilisation und Transport in eine geeignete Einrichtung
zur Abklärung}


Vgl. \prettyref{m:einschaetzung-indikation-ws-immobilisation}
\begin{itemize}
  \item \EE{\Abk{HWS}-Immobilisation} (HWS-Schiene)
  \item Bergung mit Schaufeltrage, Spineboard oder
  KED-System
  \item \EE{Lagerung} und Schienung mittels Vakuummatratze, Schaufeltrage mit
  Fixationsmöglichkeit oder Spineboard.\footnote{Welches Hilfsmittel
  eingesetzt wird ist abhängig vom Patientenzustand,
  dem Training des Teams, internen Vorschriften und
  vorhandenem Material.}
  \item \EE{Monitoring}: Achte besonders auf Störungen
  von Atmung und Bewusstsein, Schockentwicklung und
  beginnende neurologische Ausfälle.
%   \item \EE{Unfallmechanismus} erheben und dokumentieren
  \item \EE{Transportentscheidung}: Je nach
  lokalen Gegebenheiten
\end{itemize}
}




% % M %%%%%%%%%%%%%%%%%%%%%%%%%%%%%%%%%%%%%%%%%%%%%%%%%%% M %
% % 
% \ProcedurePrintDefault{}
\MassnahmenNG{Spezielle Maßnahmen}
{Rückenmarksverletzung mit Symptomen}
{\label{m:ws-trauma}}
{%
\EE{Taktik}: \Speech{\Ca{Zeitkritisch! Vitale Bedrohung!} Zügige Bergung,
Immobilisation und Transport in eine geeignete Einrichtung}


Vgl. \prettyref{m:einschaetzung-indikation-ws-immobilisation}
\begin{itemize}
  \item Bei ansprechbaren Patienten: immer
  \EE{\Abk{HWS}-Immobilisation} (HWS-Schiene)
  \item \EE{Bergung} mit Schaufeltrage, Spineboard oder
  KED-System
  \item \EE{Lagerung} und Schienung mittels
  mit Vakuummatratze, Schaufeltrage mit Fixationsmöglichkeit oder
  Spineboard.\footnote{Welches Hilfsmittel
  eingesetzt wird ist abhängig vom Patientenzustand,
  dem Training des Teams, internen Vorschriften und
  vorhandenem Material.}
  \item Bei Querschnittslähmung: Kennzeichnung der
  Sensibilitätsgrenze mit Angabe der Uhrzeit
  \item \EE{Monitoring}: Achte besonders auf Störungen
  von Atmung und Bewußtsein, sowie Schockentwicklung.
%   \item \EE{Unfallmechanismus} erheben und dokumentieren 
  \item \EE{Transportentscheidung}: Schockraum
\end{itemize}
}




% % M %%%%%%%%%%%%%%%%%%%%%%%%%%%%%%%%%%%%%%%%%%%%%%%%%%% M %
% % 
% \ProcedurePrintDefault{}
\MassnahmenNG{Spezielle
Maßnahmen}{Pneumothorax}{\label{m:pneumothorax}}
{
\EE{Taktik}: \Speech{\Ca{Zeitkritisch! Vitale Bedrohung!}
Entlastungspunktion, zügiger Transport
in eine geeignete Einrichtung}

\begin{itemize}
  \item \TxMassVitMKBes
  \begin{itemize}
    \item \EE{Lagerung}: Oberkörper hoch. Wenn stabile
    Seitenlage notwendig ist, dann Lagerung auf die
    \E{verletzte} Seite, damit die unverletzte
    Lunge  nicht unnötig belastet wird.
    \item \Z{\EE{Notarzt}: Ein Pneumothorax kann mittels
    Drainage bzw. Punktion entlastet werden (ärztliche Maßnahme).}
  \end{itemize}
  \item \EE{Verband} bei offener Thoraxverletzung: \EE{Nicht
  luftdicht verbinden}, es kann sonst ein
  Spannungspneumothorax entstehen! Steril abdecken.
  \item \EE{Transportentscheidung}: Schockraum
\end{itemize}
}




% % M %%%%%%%%%%%%%%%%%%%%%%%%%%%%%%%%%%%%%%%%%%%%%%%%%%% M %
% % 
% \ProcedurePrintDefault{}
\MassnahmenNG{Spezielle Maßnahmen}{Bauchtrauma,
offen}{\label{m:bauchtraume-offen}}{
\EE{Taktik}: \Speech{\Ca{Zeitkritisch! Vitale Bedrohung!}
Zügiger Transport in eine geeignete Einrichtung}

\begin{itemize}
  \item \TxMassVitMKBes
  \begin{itemize}
    \item \EE{Lagerung}: Bauchdecke entspannen , Knierolle
  \end{itemize}
  \item Darmteile nicht  zurückstopfen ${\rightarrow}$ sonst
  Darmriß/Darmverschluß
  \item Steril abdecken, mit Ringerlösung feucht halten
  \begin{itemize}
    \item Gefahr des Eindringens von Keimen  (Peritonitis/Sepsis!)
  \end{itemize}
  \item \EE{Transportentscheidung}: Schockraum
\end{itemize}
}




% % M %%%%%%%%%%%%%%%%%%%%%%%%%%%%%%%%%%%%%%%%%%%%%%%%%%% M %
% % 
% \ProcedurePrintDefault{}
\MassnahmenNG{Spezielle Maßnahmen}{Beckentrauma,
instabil}{\label{m:beckentrauma-instabil}}
{
\begin{itemize}
  \item \TxMassVitMKBes
  \begin{itemize}
    \item Lagerung: flach
  \end{itemize}
  \item Schockbekämpfung
  \item Wirbelsäulenschonende, situationsgerechte Rettung (z.\,B. mittels
  Schaufeltrage)
  \item Lagerung auf/Schienung mit Vakuummatratze (wenn
  Lagerung mit erhöhten Beinen notwendig ist,
  Tragentisch schräg stellen)
  \item Bei Beckenfraktur:
  \index{Beckengurt}\EE{Beckengurt}, zur Not mit
  Gürtel
  \item Transportentscheidung: Schockraum
\end{itemize}
}




% % M %%%%%%%%%%%%%%%%%%%%%%%%%%%%%%%%%%%%%%%%%%%%%%%%%%% M %
% % 
% \ProcedurePrintDefault{}
\MassnahmenNG{Spezielle
Maßnahmen}{Polytrauma}{\label{m:polytrauma}}
{%
\EE{Taktik}: \Speech{\Ca{Vitale Bedrohung! Zeitkritisch!}
Rasche Bergung, Stabilisierung
der Vitalfunktionen, Immobilisierung, rasche Versorgung von
\E{gefährlichen} Verletzungen, rascher Transport}


\begin{itemize}
  \item \TxMassVitMKBes
  \begin{itemize}
    \item Lagerung: situationsgerecht, ggfs.
    Wirbelsäulenimmobilisation, bzw. entsprechend den
    Verletzungen und der Herz-Kreislauf-Situation
%     -- Prioritäten absteigend:
%     {
%     \scriptsize
%     \begin{enumerate}
%       \item Atem- und Kreislaufstillstand $\rightarrow$ Rückenlage
%       \item Bewusstlosigkeit, nicht intubiert $\rightarrow$ stabile Seitenlage 
%       (bei Thorax- und/oder Extremitätentrauma auf verletzte Seite legen)
%       \item Atemnot bei Thoraxtrauma $\rightarrow$ Oberkörper hoch (45\textdegree)
%       \item Abdominaltrauma $\rightarrow$ Schonhaltung, Knierolle
%       \item \Abk{SHT} $\rightarrow$ Oberkörper hoch (30\textdegree)
%       \item Wirbelsäulentrauma $\rightarrow$ Flachlagerung auf Vakuummatratze
%       \item Extremitätentrauma $\rightarrow$ Extremität schienen und fixieren
%       \item Schock $\rightarrow$ Hochlagerung der Beine
%       durch schräg stellen des Tragentisches
%     \end{enumerate}
%     }
  \end{itemize}
  \item Gefährliche Verletzungen versorgen (starke
  Blutungen, Beckengurt, \ldots) 
  \item \EE{Wärmeerhalt}: Der Wärmeerhalt hat massiven Einfluss auf das
  Überleben!
%   \item Wenn einigermaßen stabil, gründlicher  Traumacheck
  \item Verletzungsmuster/Unfallmechanismus erheben $\rightarrow$ zu 
  erwartende Komplikationen einschätzen
  \item \Z{Notarzt: Schmerztherapie, Narkose, Intubation und
  Beatmung}
  \item Transportentscheidung:
  \begin{itemize}
    \item Transportentscheidung: Schockraum, mit
    Voranmeldung (Aviso) über Leitstelle
    \item Hubschrauber-Transport erwägen
  \end{itemize}
\end{itemize}
}




% % M %%%%%%%%%%%%%%%%%%%%%%%%%%%%%%%%%%%%%%%%%%%%%%%%%%% M %
% % 
% \ProcedurePrintDefault{}
\MassnahmenNG{Spezielle Maßnahmen}{Versorgung eines
Amputats}{\label{m:amputat}}
{%
\begin{itemize}
  \item Keine Reinigung des Amputats
  \item Sterile und trockene Versorgung des Amputats
  (Verband)
  \item Transport in doppelwandigem
  Amputatsbeutel (Replant-Beutel) bei  einer Temperatur von
  ca. 4\,\DegC \cite{BoehmerTaschRett6,Gorgass7,LPN3}.
  Zwischen die beiden Schichten des Beutels gibt man Wasser
  mit Eis (im Verhältnis 1:1).
  \item Das Amputat darf auf keinen Fall einfrieren -- im
  Zweifelsfall eher wärmer als kälter lagern. Auch der
  direkte Kontakt zu Wasser oder Eis muss unbedingt
  vermieden werden!
  \item \Z{Ausgeschlagene Zähne sollten in einer geeigneten
  antiseptischen Flüssigkeit transportiert werden.
  \cite{BoehmerTaschRett6,Gorgass7}}
\end{itemize}
}




% % M %%%%%%%%%%%%%%%%%%%%%%%%%%%%%%%%%%%%%%%%%%%%%%%%%%% M %
% % 
% \ProcedurePrintDefault{}
\MassnahmenNG{Spezielle Maßnahmen}{Verbrennung, leicht}{\label{m:verbrennung-leicht}}
{    
\begin{itemize}
  \item Selbstschutz!
  \item Kleiderbrand löschen
  \item Vitale Bedrohung einschätzen. \TxBeiVit \TxMassVitMK
  \item Kleidung und Schmuck entfernen (sofern nicht  eingeschmolzen!)
  \item Nur wenn \E{unmittelbar nach dem Unfall\footnote{Bereits zwei
  Minuten nach Verbrennungsbeginn und damit bei Eintreffen
  des Rettungsdienstes ist ein
  positiver Effekt der Kühlung nicht mehr zu erwarten. Bei
  mehreren Minuten zurückliegenden Verbrennungen ist eine
  Kühlung nicht mehr sinnvoll.
  \cite{Asanger2010,AsboeLehrmeinungRd2011-21}}}
  damit begonnen werden kann, für max. 10\Min* mit handwarmen
  Wasser die betroffenen Körperregionen spülen;
  
  \EE{Bei Frösteln des Patienten früher abzubrechen: Gefahr der Unterkühlung!}
  \item \EE{Abdeckung}: Sterile, nicht verklebende
  Wundauflagen; evtl. metalliserte Wundauflagen
  \cite{Helfen2008,RedelsteinerHRNS1}; locker anlegen
  % DISCUSS INHALT: Water-Jel, Stellungsnahme
  %         \item Wenn vorhanden: Water-Jel
  \item Ggfs. Schockbekämpfung
  %   \item Ggfs. Notarzt-Anforderung
\end{itemize}
}




% % M %%%%%%%%%%%%%%%%%%%%%%%%%%%%%%%%%%%%%%%%%%%%%%%%%%% M %
% % 
% \ProcedurePrintDefault{}
\MassnahmenNG{Spezielle Maßnahmen}{Verbrennung, schwer}{\label{m:verbrennung-schwer}}
{%
\begin{itemize}
  \item \TxMassVitMK
  \item Kleidung und Schmuck entfernen (sofern nicht  eingeschmolzen!)
  \item Besonders auf \EE{Wärmeerhalt} achten!
  \item Luftzug vermeiden
  \item \Z{\E{Keine} Kühlung\footnote{Bei großflächigen
  Verbrennungen sollte der Patient nur primär abgelöscht werden. Hier soll
  keine Kühlung der betroffenen Hautregionen erfolgen.
  \cite{Asanger2010}}}
%   \item Mind. 10 min. mit mäßig kaltem Wasser die betroffenen Körperregionen spülen  --
%   \EE{NICHT unterkühlen!} 
%   OBSOLETE: Großflächige Verbrennungen (Erwachsene
% > 20% KOF, Kinder > 10% KOF, Säuglinge > 5%) sollten nur primär
% abgelöscht und keimfrei abgedeckt werden. Hier erfolgt keine Kühlung
% der betroffenen Regionen.
  \item \EE{Abdeckung}: Sterile, nicht verklebende
  Wundauflagen, evtl. metalliserte Wundauflagen
  \cite{Helfen2008,RedelsteinerHRNS1}.\ACh{GABS}{2010-08-06}{Metallisierte
  Verbände sind wünschenswert, werden zB in Helfen2008 
  empfohlen.}\ZFootnote{Der Nutzen  von kühlenden
  Spezialverbänden, z.\,B. Waterjel{\texttrademark} ist
  aufgrund  der Unterkühlungs-Gefahr und einer  möglichen
  Durchblutungsstörung der Wunde
  (Gefahr von Wundheilungsstörungen im weiteren Verlauf)
  umstritten.} Locker anlegen. Bei großflächigen
  Verbrennungen können,  wenn das sterile Material nicht
  ausreicht,  zur Not auch saubere Leintücher o.\,ä. 
  verwendet werden \cite{emergency9,Asanger2010}.
  % TODO INHALT: Stellungsnahme zu Water-Jel
  %         \item Wenn vorhanden: Water-Jel
  %         \footnote{Streitfrage: Water-Jel: Umstritten}
  \item Schockbekämpfung
  \item \Z{NKV+: Infusionstherapie mit \E{kristalloiden} Lösungen; kolloidale
  Lösungen sind \E{kontraindiziert}.}
\end{itemize}
}




% % M %%%%%%%%%%%%%%%%%%%%%%%%%%%%%%%%%%%%%%%%%%%%%%%%%%% M %
% % 
% \ProcedurePrintDefault{}
\MassnahmenNG{Spezielle
Maßnahmen}{Inhalationstrauma}{\label{m:inhalationstrauma}}{

\begin{itemize}
  \item Vitale Bedrohung einschätzen. \par
  \TxBeiVit \TxMassVitMK
  \item Lagerung: Oberkörper hoch
  \item \ce{O2}-Gabe gemäß \prettyref{m:sauerstoffberieselung}
  \item Engmaschiges Monitoring
\end{itemize}

% \Customized{
% \begin{description}
%   \item[\CustAsbFlo] Notfallsanitäter geben bei einer Rauch- oder
%   Reizgasinhalation ein inhalatives Kortikoid (z.\,B.
%   Pulmicort{\texttrademark}) laut Algorithmus der
%   Medikamentenliste~1.
% \end{description}
% }
}




% % M %%%%%%%%%%%%%%%%%%%%%%%%%%%%%%%%%%%%%%%%%%%%%%%%%%% M %
% % 
% \ProcedurePrintDefault{}
\MassnahmenNG{Spezielle Maßnahmen}{Verletzungen mit chemischen
Substanzen}{\label{m:chemische-verletzungen}}{

\begin{enumerate}
  \item Auf Selbsschutz achten!
  \begin{itemize}
    \item Patienten ggfs. retten, aber nur wenn es der
    Eigenschutz zulässt.
  \end{itemize}
  \item Gefahr erkennen: \EE{Welcher Stoff?}  Expertenrat einholen!
  \item Wenn erforderlich: Vorgehen analog zu Abschnitt
  \enquote{Gefahrengutunfall}, GAS-Regel
  (\prettyref{AASS-topic:gefahrengutunfall})
  \item Weitere Opfer ausschließen (lassen)
  \item Sicherung der Vitalfunktionen, Beurteilung ob der
  Patient vital bedroht ist.
  \item \TxBeiVit \TxMassVitMKBes
  \begin{itemize}
    \item Lagerung: situationsgerecht.
  \end{itemize}
  \item Entfernung kontaminierter Kleidung
  \item Spülung der betroffenen Hautpartien mit Wasser  (richtig abrinnen
  lassen)
  \item Auge: Spülung mit Wasser / Ringerlösung (nach außen abrinnen
  lassen!),  Kontaktlinsen entfernen
  \item Transportentscheidung: Ab einer Ausdehnung von
  15\PC* (Erwachsener) bzw. 10\PC* (Kind) Spezialabteilung
  für Verbrennungsverletzungen. Bei Augenverletzungen Abt.
  für Augenheilkunde.
\end{enumerate}
}




% % M %%%%%%%%%%%%%%%%%%%%%%%%%%%%%%%%%%%%%%%%%%%%%%%%%%% M %
% % 
% \ProcedurePrintDefault{}
\MassnahmenNG{Spezielle
Maßnahmen}{Erfrierungen}{\label{m:erfrierung}}
{%
\begin{itemize}
  \item Beengende Kleidungsstücke lockern (Einschnürung erhöht
  Risiko einer Erfrierung)
  \item Keimfreier, gepolsterter Verband
  \item Warme, gezuckerte Getränke -- KEIN Alkohol
  \item Körperstamm wärmen, NIE den erfrorenen Körperteil abreiben!!
\end{itemize}
}




% % M %%%%%%%%%%%%%%%%%%%%%%%%%%%%%%%%%%%%%%%%%%%%%%%%%%% M %
% % 
% \ProcedurePrintDefault{}
\MassnahmenNG{Spezielle
Maßnahmen}{Stromunfälle}{\label{m:stromunfall}}
{
\begin{itemize}
  \item \EE{Eigenschutz!} Bevor der Patienten berührt wird,
  muss man sich vergewissern, dass der Patient
  keinen Kontakt mehr mit dem Stromkreis hat! \EE{Strom
  abschalten (lassen)!}
  \begin{itemize}
    \item Gefahrenbreich beachten:
    \prettyref{AASS-topic:gefahrenbereich-strom}
    \item Sicherheitsabstand
    (\prettyref{AASS-topic:gefahrenbereich-strom})
    \item Schrittspannung und Spannungskegel
    beachten (\prettyref{AASS-topic:gefahrenbereich-strom})
    \item Bahnanlagen sperren/sichern lassen
    (\prettyref{AASS-topic:gefahrenbereich-strom})
    \item Schalter ausschalten, (Haupt-)Sicherung
    auslösen, Starkstrom/Hochspannung durch Fachmann  abschalten bzw.
    erden lassen
    \item Strom gegen Wiedereinschalten sichern!
    \item Patient mit isolierendem Material aus dem Stromkreis befreien

  \end{itemize}
  \item Beurteilung der vitalen Bedrohung.
  \par \TxBeiVit
  \TxMassVitMK
  \item \EE{Monitoring und
  Reanimationsbereitschaft}: auf evtl.
  Herzrhythmusstörungen vorbereitet sein
  (bis zu 24\,h Latenz)
  \item Genaue Patientenuntersuchung ${\rightarrow}$
  Eintritts- und Austrittsmarke
  \item Hospitalisierung auch bei symptomlosen
  Patienten zur Überwachung
  \item Stromverbrennungen wie Verbrennungen versorgen
\end{itemize}
}





% \appendix
% {\WideParBegin}
% \part{Appendix}
% {\WideParEnd}
%\setcounter{page}{0}

% \pagenumbering[Appendix]{bychapter}


% \chapter{Verzeichnisse}



%%%%%%%%%%%%%%%%%%%%%%%%%%%%%%%%%%%%%%%%%%%%%%%%%%%%%%%%%%%
%%% Kommentare
\printnotes

%%%%%%%%%%%%%%%%%%%%%%%%%%%%%%%%%%%%%%%%%%%%%%%%%%%%%%%%%%%
%%% Literaturverzeichnis
% \opt{STATUS-WITH-BIBLIO}{{\loadgeometry{FullPage}\twocolumn\tiny\setlength{\parskip}{0pt}\cleardoublepage\printbibliography\restoregeometry}\onecolumn}
% \opt{STATUS-WITH-BIBLIO}{{\setlength{\parskip}{0pt}\bibliography{Literatur-AASS}}}
{\loadgeometry{FullPage}\twocolumn\setlength{\parskip}{0pt}\cleardoublepage\printbibliography\restoregeometry\twocolumn}







\pagestyle{empty}\loadgeometry{WholePage}
\small\flushleft

\chapter{Übersichten und Tabellen}


\Layout{57}{}{}{}{
\begin{tabular}{m{2.55cm}m{2.5cm}m{2.55cm}m{2.5cm}}
\toprule\addlinespace
\multicolumn{2}{m{5.1cm}}{\centering \TabHeading{Allgemeine
Kennzeichnung} (Sammeltransporte)} &
\multicolumn{2}{m{5.1cm}}{\centering \TabHeading{Spezielle
Kennzeichnung}} \\\addlinespace\midrule\addlinespace
\includegraphics[width=\linewidth]{./images/khd/warntafel-allgemein.png}
 &
 &
\includegraphics[width=\linewidth]{./images/khd/warntafel-un-60-1710.png}
 &
Gefahrnummer\par (\T{Kemler-Nummer})

Stoffnummer\par (\T{UN-Nummer})
\\\addlinespace\bottomrule
\end{tabular}
}{Kennzeichnung durch Gefahrentafel.}{}




\Layout{57}{}{}{}{
\begin{tabulary}{\linewidth}{cL}
\addlinespace\toprule\addlinespace
 2 & Entweichen von Gas durch Druck oder chemische Reaktion
\\\addlinespace
 3 & Entzündbarkeit von Flüssigkeiten (Dämpfen) und Gasen oder 
selbsterhitzungsfähiger flüssige
Stoffe (z.B. Benzin, Diesel)
\\\addlinespace
 4 & Entzündbarkeit von festen Stoffen oder
selbsterhitzungsfähiger fester Stoffen
\\\addlinespace
 5 & Oxidierende (brandfördernde) Wirkung
\\\addlinespace
 6 & Giftigkeit oder Ansteckungsgefahr
\\\addlinespace
 7 & Radioaktivität
\\\addlinespace
 8 & \"Atzwirkung
\\\addlinespace
 9 & Gefahr einer spontan heftigen Reaktion
\\\addlinespace
0 & Keine weitere Gefahr
\\\addlinespace
 {\sffamily X} & Reagiert gefährlich mit Wasser
\\\addlinespace\bottomrule
\end{tabulary}
}{Bedeutung der
Ziffern der Gefahrennummer\label{tab:gefahrennummer}}{}





\Layout{57}{}{}{}{
\begin{tabulary}{\linewidth}{cL}
\addlinespace\toprule\addlinespace
22 & Tiefgekühltes Gas
\\\addlinespace
X323 & Entzündbarer flüssiger Stoff, der mit Wasser gefährlich reagiert, wobei entzündbare Gase entweichen
\\\addlinespace
X333 & Selbstentzündliche Flüssigkeit, die mit Wasser gefährlich reagiert
\\\addlinespace
X423 & Entzündbarer fester Stoff, der mit Wasser gefährlich reagiert, wobei brennbare Gase entweichen
\\\addlinespace
44 & Entzündbarer fester Stoff, der sich bei erhöhter Temperatur in Geschmolzenem Zustand befindet
\\\addlinespace
539 & Entzündbares organisches Peroxid
\\\addlinespace
90  & Verschiedene gefährlich Stoffe
\\\addlinespace\bottomrule
\end{tabulary}
}
{Spezielle Ziffernkombinationen bei der
Gefahrennummer\label{tab:gefahrennummer-spezielle-ziffernkombinationen}}
{}










\Layout{57}{}
{}
{}
{
\begin{tabulary}{\linewidth}{cLcL}
\toprule\addlinespace
0048 & Sprengkörper & 1223 & Kerosin
\\\addlinespace
0333 & Feuerwerkskörper & 1789 & Salzsäure
\\\addlinespace
1005 & Ammoniak & 1830 & Schwefelsäure
\\\addlinespace
1011 & Butan & 1962 & Ethylen
\\\addlinespace
1013 & Kohlendioxid & 1972 & Methan oder Erdgas (tiefgekühlt, flüssig)
\\\addlinespace
1017 & Chlor & 1977 & Stickstoff
\\\addlinespace
1052 & Flußsäure & 1978 & Propan
\\\addlinespace
1090 & Aceton & 3295 & Kohlenwasserstoff, flüssig
\\\addlinespace
1203 & Benzin & 3312 & Gas, tiefgekühlt, flüssig
\\\addlinespace\bottomrule
\end{tabulary}
}
{Ausgewählte Stoffnumern}
{}








\begin{WidePar}
\FormatTable\scriptsize

\begin{tabularx}{\linewidth}{XXXXXXXXX}
\toprule\addlinespace
\multicolumn{9}{c}{{\normalsize\TabSub{Gefahrgutsymbole}}}
\\\addlinespace \multicolumn{9}{c}{\Z{\footnotesize nach Verordnung
(EG) Nr. 1272/2008 gültig ab 1.12.2010 \cite{EG-1272-2008-De}}}
\\\addlinespace\midrule\addlinespace

\includegraphics[width=\linewidth]{./images/GEF/gefahrensymbole/ghs01_explodingbomb.png}
&

\includegraphics[width=\linewidth]{./images/GEF/gefahrensymbole/ghs02_flame.png}
&

\includegraphics[width=\linewidth]{./images/GEF/gefahrensymbole/ghs03_flameovercircle.png}
&

\includegraphics[width=\linewidth]{./images/GEF/gefahrensymbole/ghs04_gascylinder.png}
&

\includegraphics[width=\linewidth]{./images/GEF/gefahrensymbole/ghs05_corrosion.png}
&

\includegraphics[width=\linewidth]{./images/GEF/gefahrensymbole/ghs06_skullandcrossbones.png}
&

\includegraphics[width=\linewidth]{./images/GEF/gefahrensymbole/ghs07_exclamationmark.png}
&

\includegraphics[width=\linewidth]{./images/GEF/gefahrensymbole/ghs08_healthhazard.png}
&

\includegraphics[width=\linewidth]{./images/GEF/gefahrensymbole/ghs09_environment.png}
\\\addlinespace
Gefahr:\newline
Explosions\-gefährlich
&
Gefahr:\newline
Leicht-/Hoch\-entzündlich 
&
Gefahr:\newline
Brandfördernd 
&
Achtung:\newline
Komprimierte Gase 
&
Gefahr:\newline
Ätzend 
&
Gefahr:\newline
Giftig/ Sehr giftig 
&
Achtung:\newline
Gesundheits\-gefährdend 
&
Gefahr:\newline
Gesundheits\-schädlich  
&
Warnung:\newline
Umwelt\-gefährdend
\\\addlinespace
\Z{\tiny GHS01}
&
\Z{\tiny GHS02}
&
\Z{\tiny GHS03}
&
\Z{\tiny GHS04}
&
\Z{\tiny GHS05}
&
\Z{\tiny GHS06}
&
\Z{\tiny GHS07}
&
\Z{\tiny GHS08}
&
\Z{\tiny GHS09}
\\\addlinespace\midrule\addlinespace
\multicolumn{9}{c}{Diese GHS-Symbole werden mit
Gefahrenhinweisen und Sicherheitshinweisen ergänzt!}
\\\addlinespace\toprule\addlinespace
\multicolumn{9}{c}{{\normalsize \TabSub{Alte Zeichen und
Gefahrenbezeichnungen} (Bis Mitte 2017)}}
\\\addlinespace\midrule\addlinespace
\includegraphics[width=\linewidth]{./images/GEF/gefahrensymbole/gefahr5.png}
&
\includegraphics[width=\linewidth]{./images/GEF/gefahrensymbole/gefahr1.png}
&
\includegraphics[width=\linewidth]{./images/GEF/gefahrensymbole/gefahr2.png}
&
&
\includegraphics[width=\linewidth]{./images/GEF/gefahrensymbole/gefahr3.png}
&
\includegraphics[width=\linewidth]{./images/GEF/gefahrensymbole/gefahr6.png}
&

&

&
\includegraphics[width=\linewidth]{./images/GEF/gefahrensymbole/gefahr4.png}
\\\addlinespace
\multicolumn{1}{c}{E} 
&
\multicolumn{1}{c}{F, F+} 
&
\multicolumn{1}{c}{O} 
&
&
\multicolumn{1}{c}{C} 
&
\multicolumn{1}{c}{T, T+} 
&

&

&
\multicolumn{1}{c}{N} 
\\\addlinespace
Explosions\-gefährlich 
&
Leicht-/Hoch\-entzündlich 
&
Brandfördernd 
&
keine Entsprechung 
&
Ätzend 
&
Giftig/ Sehr giftig
&
keine Entsprechung
&
keine Entsprechung 
&
Umwelt\-gefährlich
\\\addlinespace\bottomrule
\end{tabularx}
\end{WidePar}







\Layout{57}{}
{}
{}
{% \begin{WidePar}
{\footnotesize
\begin{tabular}{m{0.28\linewidth}m{0.15\linewidth}m{0.28\linewidth}m{0.15\linewidth}}
\toprule\addlinespace
% \TabHeading{Gefahrenklasse} 
% &
% \TabHeading{Symbol}
% 
% (Auswahl)
% \\\addlinespace\midrule\addlinespace
\TabSub{Klasse 1}

Explosive Stoffe und Gegenstände mit Explosivstoff

Unterklassen:\newline
1.1, 1.2, 1.3, 1.4, 1.5, 1.6 
&
~\newline
\includegraphics[width=1cm]{./images/GEF/gefahrgut/gefahrgut_1.png}
\\\addlinespace\midrule\addlinespace
\TabSub{Klasse~2.1}

Entzündbare Gase 
&
~\newline
\includegraphics[width=1cm]{./images/GEF/gefahrgut/gefahrgut_21.png}
&
\TabSub{Klasse~2.2}

Nicht entzündbare,\newline
nicht giftige Gase 
&
~\newline
\includegraphics[width=1cm]{./images/GEF/gefahrgut/gefahrgut_22.png}
\\
\TabSub{Klasse~2.3}

Giftige Gase 
&
~\newline
\includegraphics[width=1cm]{./images/GEF/gefahrgut/gefahrgut_23.png}
 \\\addlinespace\midrule\addlinespace
\TabSub{Klasse~3}

Entzündbare\newline
flüssige Stoffe 
&
~\newline
\includegraphics[width=1cm]{./images/GEF/gefahrgut/gefahrgut_3.png}
\\\addlinespace\midrule\addlinespace
\TabSub{Klasse~4.1}

Entzündbare feste Stoffe, selbstzersetzliche Stoffe und
desensibilisierte explosive feste Stoffe 
&
~\newline
\includegraphics[width=1cm]{./images/GEF/gefahrgut/gefahrgut_41.png}
&
\TabSub{Klasse~4.2}

Selbstentzündliche Stoffe 
&
~\newline
\includegraphics[width=1cm]{./images/GEF/gefahrgut/gefahrgut_42.png}
\\
\TabSub{Klasse~4.3}

Stoffe, die in Berührung mit Wasser entzündliche Gase entwickeln 
&
~\newline
\includegraphics[width=1cm]{./images/GEF/gefahrgut/gefahrgut_43.png}
\\\addlinespace\midrule\addlinespace
\TabSub{Klasse~5.1}

Entzündend\newline
(oxidierend)\newline
wirkende Stoffe 
&
\includegraphics[width=1cm]{./images/GEF/gefahrgut/gefahrgut_51.png}
&
\TabSub{Klasse~5.2}

Organische Peroxide
&
~\newline
\includegraphics[width=1cm]{./images/GEF/gefahrgut/gefahrgut_52.png}
\\\addlinespace\midrule\addlinespace
\TabSub{Klasse~6.1}

Giftige Stoffe 
&
~\newline
\includegraphics[width=1cm]{./images/GEF/gefahrgut/gefahrgut_61.png}
&
\TabSub{Klasse~6.2}

Ansteckungs\-gefährliche Stoffe 
&
~\newline
\includegraphics[width=1cm]{./images/GEF/gefahrgut/gefahrgut_62.png}
\\\addlinespace\midrule\addlinespace
\TabSub{Klasse~7}

Radioaktive Stoffe 
&
\includegraphics[width=1cm]{./images/GEF/gefahrgut/gefahrgut_7.png}
~
\includegraphics[width=1cm]{./images/GEF/gefahrgut/gefahrgut_7_fissile.png}
\\\addlinespace\midrule\addlinespace
\TabSub{Klasse~8}

Ätzende Stoffe 
&
\includegraphics[width=1cm]{./images/GEF/gefahrgut/gefahrgut_8.png}
\\\addlinespace\midrule\addlinespace
\TabSub{Klasse~9}

Verschiedene gefährliche Stoffe und Gegenstände &


\includegraphics[width=1cm]{./images/GEF/gefahrgut/gefahrgut_9.png}
\\\addlinespace\midrule\addlinespace
\end{tabular}
% \end{WidePar}
}
}
{\EE{Gefahrenklassen} und ausgewählte Symbole
(Gefahrenzettel), nach \citep{ADR2011}} {}


\begin{WidePar}
\begin{center}
\vfill
\includegraphics[height=.9\textheight]{./images/emhofer-josef-aass/BLS}
\vfill
\Captionof{figure}{Basic Life Support mit AED\index{Basic Life Support}\index{BLS|see{Basic Life Support}}}
\end{center}
\end{WidePar}

\begin{WidePar}
\begin{center}
\includegraphics[height=.9\textheight]{./images/emhofer-josef-aass/ALS}
\vfill
\Captionof{figure}{Advanced Life Support mit AED\index{Advanced Life Support}\index{ALS|see{Advanced Life Support}}}
\end{center}
\end{WidePar}

\begin{WidePar}
\begin{center}
\includegraphics[height=.9\textheight]{./images/emhofer-josef-aass/PLS}
\vfill
\Captionof{figure}{Pediatric Life Support mit AED\index{Pediatric Life Support}\index{PLS|see{Pediatric Life Support}}}
\end{center}
\end{WidePar}


\Layout{57}{}{}{}{
\begin{tabularx}{\linewidth}{llXXXXXXX}%
\addlinespace\toprule\addlinespace
\noindent
&
&
\noindent\TabHeading{NG}
&
\noindent\TabHeading{Sg}
&
\noindent\TabHeading{KK}
&
\noindent\TabHeading{SK}
&
\noindent\TabHeading{Ju}
&
\noindent\TabHeading{Erw \Male}
&
\noindent\TabHeading{Erw \Female}
\\\addlinespace\midrule\addlinespace
\TabSub{RR\textsubscript{sys}}
&
[\MetricUnit{mmHg}]
&
75\nocite{Soehnke200607}\nocite{Lissauer2007}
&
\VonBis{80}{90}\nocite{Soehnke200607}
&
95\nocite{Lissauer2007}
&
\VonBis{100}{110}\nocite{Soehnke200607}\nocite{Lissauer2007}
&
120\nocite{Soehnke200607}
&
\VonBis{100}{140}\nocite{Soehnke200607}
&
$\hookleftarrow$
\\\addlinespace
\TabSub{HF} 
&
[\PerMin]
&
\VonBis{140}{180}
&
\VonBis{110}{160}\nocite{Lissauer2007}
&
\VonBis{95}{140}\nocite{Lissauer2007}
&
\VonBis{80}{120}\nocite{Lissauer2007}
&
\VonBis{60}{100}\nocite{Lissauer2007}
&
\VonBis{60}{100}
&
$\hookleftarrow$
\\\addlinespace\midrule\addlinespace
\TabSub{AF} 
&
[\PerMin]
&
\VonBis{30}{50}\nocite{Lissauer2007}
&
\VonBis{20}{30}\nocite{Lissauer2007}
&
\VonBis{20}{30}\nocite{Lissauer2007}
&
\VonBis{15}{20}\nocite{Lissauer2007}
&
\VonBis{14}{20}
&
\VonBis{12}{16}
&
$\hookleftarrow$
% \\\addlinespace
% \TabSub{AZV} 
% &
% [\nicefrac{\Ml}{\Kg}]
% &
% 
% &
% 
% &
% 
% &
% 
% &
% 
% &
% 
% &
% 
\\\addlinespace
\TabSub{AZV}\footnote{Bezogen auf Normalgewicht, \E{nicht}
auf Ist-Gewicht!} 
&
[\Ml]
&
\VonBis{20}{30}\nocite{LutomskyLeitRett3}
&
40\,--\,55\,--\,80\nocite{LutomskyLeitRett3}
&
\VonBis{80}{180}\nocite{LutomskyLeitRett3}
&
\VonBis{240}{350}\nocite{LutomskyLeitRett3}
&
500\nocite{LutomskyLeitRett3}
&
800\nocite{LutomskyLeitRett3}
&
700\nocite{LutomskyLeitRett3}
% \\\addlinespace
% \TabSub{AMV} 
% &
% [\nicefrac{\Liter}{\Min}]
% &
% \VonBis{600}{1500}
% &
% 
% &
% 
% &
% 
% &
% 
% &
% 
% &
% 
\\\addlinespace\midrule\addlinespace
\TabSub{BZ}\footnote{nüchtern}
&
[\MgDl]
&
$\hookrightarrow$
&
$\hookrightarrow$
&
$\hookrightarrow$
&
$\hookrightarrow$
&
$\hookrightarrow$
&
\VonBis{80}{100}\nocite{WHO2006Diabetes}
&
$\hookleftarrow$
\\\addlinespace
\TabSub{Sp\ce{O2}}
&
[\PC]
&
$\hookrightarrow$
&
$\hookrightarrow$
&
$\hookrightarrow$
&
$\hookrightarrow$
&
$\hookrightarrow$
&
\VonBis{95}{100}
&
$\hookleftarrow$
\\\addlinespace
\TabSub{Temp}\footnote{Körperkerntemperatur}
&
[\DegC]
&
$\hookrightarrow$
&
$\hookrightarrow$
&
$\hookrightarrow$
&
$\hookrightarrow$
&
$\hookrightarrow$
&
37
&
$\hookleftarrow$
\\\addlinespace\bottomrule
\end{tabularx}
}{Übersicht Normalwerte}{}

\Layout{57}
{}
{}
{}
{\small{\begin{tabularx}{\linewidth}{XXXX}
\addlinespace\toprule\addlinespace
\noindent\TabHeading{Kriterium} & \noindent\TabHeading{0 Punkte} & \noindent\TabHeading{1 Punkt} & \noindent\TabHeading{2 Punkte}
\\\addlinespace\midrule\addlinespace
\noindent\TabSub{Herzfrequenz} & unter 60\PerMin* & \Range{60}{99}\PerMin* & über 100\PerMin*  
\\\addlinespace
\noindent\TabSub{Atemanstrengung} & keine & unregelmäßig, flach &  regelmäßig, Kind schreit
\\\addlinespace
\noindent\TabSub{Reflexe} & keine & Grimassieren &  kräftiges Schreien
\\\addlinespace
\noindent\TabSub{Muskeltonus} & schlaff & leichte Beugung der Extremitäten & aktive Bewegung der Extremitäten 
\\\addlinespace
\noindent\TabSub{(Haut-) Farbe} & blau, blass & Stamm rosig, Extremitäten blau & gesamter Körper rosig 
\\\addlinespace\bottomrule\addlinespace
\end{tabularx}
}}
{Apgar-Score}
{}






SO LONG.......... AND THANKS FOR ALL THE FISH
\MultiColListEnd{lomassnahmen}
\MultiColListEnd{lofazit}
\MultiColListEnd{toc}

\end{document}